% Capítulo: La verdad de Dios (generado desde 13_Verdad/13_Verdad_capitulo.tex)
% Texto ancla: Jeremías 10:10. Clasificación: Relación; incomunicable como fuente, comunicable como veracidad.

\section{Premisa: toda verdad es verdad de Dios}

Para una cosmovisión bíblica, la verdad no es un constructo social ni una abstracción filosófica, sino la correspondencia exacta entre lo que creemos y la realidad tal como Dios la concibió, la creó y la sostiene. Si Dios es la fuente inagotable de la verdad, todo lo que es verdadero en el universo procede de él: lo que estudia la física, lo que ordena la política y la economía, lo que sostiene la moral y lo que define al ser humano.

\precision{la fórmula ``toda verdad es verdad de Dios'' condensa una idea que la tradición cristiana remonta a Agustín: dondequiera que el cristiano encuentre la verdad, debe reconocer que pertenece a su Señor (\emph{De doctrina christiana} II.18.28).}

La pregunta que recorre este capítulo la formula Pilato durante el juicio de Jesús:

\begin{verso}{Juan 18:37–38}
Le dijo entonces Pilato: ¿Luego, eres tú rey? Respondió Jesús: Tú dices que yo soy rey. Yo para esto he nacido, y para esto he venido al mundo, para dar testimonio a la verdad. Todo aquel que es de la verdad, oye mi voz. Le dijo Pilato: ¿Qué es la verdad?
\end{verso}

Pilato pregunta \emph{qué} es la verdad (\grc{τί ἐστιν ἀλήθεια;}) ante aquel que dirá \emph{quién} es la verdad (Jn 14:6). El estudio recorre ese mismo camino en cuatro etapas: qué es la verdad y de dónde le viene su firmeza, en qué esferas de la creación se la descubre, cómo choca con la percepción humana y cómo se revela, al fin, en una Persona.

El vínculo con el verso que articula esta obra es inmediato. Jeremías escribe: ``Mas alábese en esto el que se hubiere de alabar: en entenderme y conocerme, que yo soy \SENOR'' (Jer 9:24 [9:23]); y pocos versos más adelante el mismo libro declara quién es ese Dios que puede ser conocido: ``Mas \SENOR{} es el Dios verdadero'' (Jer 10:10).

%===================================================================
\section{Texto ancla: Jeremías 10:10}\label{sec:jer10}

\subsection{Ídolos que no hablan}
Jeremías 10:1--16 es una sátira contra los ídolos de las naciones, y el profeta describe con ironía cómo se fabrican:

\begin{verso}{Jeremías 10:3–5}
Porque las costumbres de los pueblos son vanidad; porque leño del bosque cortaron, obra de manos de artífice con buril. Con plata y oro lo adornan; con clavos y martillo lo afirman para que no se mueva. Derechos están como palmera, y no hablan; son llevados, porque no pueden andar. No tengáis temor de ellos, porque ni pueden hacer mal, ni para hacer bien tienen poder.
\end{verso}

El ídolo necesita clavos para mantenerse en pie. Su estabilidad es prestada: alguien lo sostiene, alguien lo carga, alguien lo viste (10:9). No habla ni anda, y no puede hacer mal ni bien. Por eso Jeremías lo llama ``vanidad'' (\heb{הֶבֶל}, \emph{hével}), la palabra del vapor y del soplo que tanto emplea Eclesiastés.

\precision{el contraste entre la ``firmeza clavada'' del ídolo y la firmeza de Dios (\emph{emet}) es una lectura propia de esta serie, sugerida por el verbo \heb{יְחַזְּקוּם} (``lo afirman'', 10:4) y por la raíz de \emph{emet}.}

\subsection{El Dios verdadero}
En medio de esa descripción, el poema irrumpe con una confesión:

\hebverso{וַיהוָה אֱלֹהִים אֱמֶת הוּא־אֱלֹהִים חַיִּים וּמֶלֶךְ עוֹלָם \\ מִקִּצְפּוֹ תִּרְעַשׁ הָאָרֶץ וְלֹא־יָכִלוּ גוֹיִם זַעְמוֹ}
\begin{verso}{Jeremías 10:10}
Mas \SENOR{} es el Dios verdadero; él es Dios vivo y Rey eterno; a su ira tiembla la tierra, y las naciones no pueden sufrir su indignación.
\end{verso}

\begin{filologia}{Jer 10:10}
\heb{אֱלֹהִים אֱמֶת} \emph{Elohim emet} & \heb{אמן} & ``Dios verdad'': el sustantivo \emph{emet} se pone junto a \emph{Elohim} para calificarlo. No dice que Dios \emph{tenga} verdad, sino que \emph{es} verdad. \\
\heb{אֱלֹהִים חַיִּים} \emph{Elohim jayyim} & \heb{חיה} & ``Dios vivo'': frente al ídolo, en el que ``no hay espíritu'' (10:14). \\
\heb{מֶלֶךְ עוֹלָם} \emph{mélej olam} & \heb{מלך} & ``Rey eterno'': al ídolo lo llevan (10:5); el \SENOR{} gobierna. \\
\heb{כִּי שֶׁקֶר נִסְכּוֹ} \emph{ki shéqer nisko} & \heb{שׁקר} & ``porque mentira es su imagen fundida'' (10:14). El ídolo es \emph{shéqer}, lo contrario de \emph{emet}. \\
\end{filologia}

Todo el capítulo enfrenta dos realidades: el Dios que \emph{es} verdad y la imagen que \emph{es} mentira. ``Todo hombre se embrutece, y le falta ciencia; se avergüenza de su ídolo todo fundidor, porque mentirosa es su obra de fundición, y no hay espíritu en ella'' (Jer 10:14). En Jeremías, la mentira no es ante todo una frase falsa, sino un dios falso; y Pablo empleará el mismo lenguaje (Ro 1:25).

\subsection{Dos columnas que se alternan}
El poema avanza como un juicio a dos voces, en el que cada descripción de los ídolos recibe su réplica sobre el Dios verdadero:

\begin{center}\small\renewcommand{\arraystretch}{1.2}
\begin{tabular}{@{}L{3.0cm}L{4.4cm}L{3.0cm}L{4.0cm}@{}}
\toprule
\textbf{Los ídolos} & & \textbf{El \SENOR} & \\ \midrule
vv. 2--5 & leño, clavos, no habla & vv. 6--7 & ``no hay semejante a ti'' \\
vv. 8--9 & ``enseñanza de vanidades'' & v. 10 & Dios verdad, vivo, Rey eterno \\
v. 11 & los dioses desaparecerán & vv. 12--13 & hizo la tierra con su poder \\
vv. 14--15 & mentira, sin espíritu & v. 16 & ``la porción de Jacob'' \\
\bottomrule
\end{tabular}\par
{\footnotesize\color{gris}\textit{Tabla 1.} La alternancia de Jeremías 10:1--16. \textit{Nota.} Elaboración propia.}
\end{center}

El versículo 11 se aparta del resto porque está escrito en arameo, la lengua internacional del imperio:

\hebverso{אֱלָהַיָּא דִּי־שְׁמַיָּא וְאַרְקָא לָא עֲבַדוּ}

\noindent ``los dioses que no hicieron los cielos ni la tierra''. Es la frase que Israel debía pronunciar ante las naciones en su propio idioma.

\precision{los comentarios describen la estructura alternada del poema y el carácter arameo del v. 11 (Lundbom, 1999).}

\precision{(textual) la Septuaginta de Jeremías es alrededor de una séptima parte más breve que el texto masorético. En el capítulo 10 omite los vv. 6--8 y el v. 10, y coloca el v. 9 después del 5a. Qumrán conserva ambas formas: 4QJer\textsuperscript{b} (4Q71) coincide con la forma corta y el orden de la LXX, mientras que 4QJer\textsuperscript{a} (4Q70) sigue la forma larga del texto masorético (Tov, 2012). La RVR1960 sigue el texto masorético, que incluye la confesión del v. 10. Su contenido, en todo caso, está presente en el resto de la Escritura: ``Dios vivo'' (Dt 5:26; Sal 42:2) y ``Dios de verdad'' (Sal 31:5; Is 65:16).}

%===================================================================
\section{La palabra: la raíz de la firmeza}

\subsection{La raíz \emph{amán} en lo concreto}
El hebreo bíblico dice ``verdad'' con \emph{emet} (\heb{אֱמֶת}), palabra formada sobre la raíz \heb{א־מ־ן} (\emph{amán}), de la que procede también ``amén''. La raíz evoca una realidad física y estructural: sostener, ser firme, edificar sobre cimientos sólidos. Sus usos concretos permiten verlo:

\begin{center}\small\renewcommand{\arraystretch}{1.18}
\begin{tabular}{@{}L{3.3cm}L{7.4cm}L{3.6cm}@{}}
\toprule
\textbf{Forma} & \textbf{Uso concreto} & \textbf{Texto} \\ \midrule
\heb{הָאֹמֵן} \emph{ha'omén} & ``la que cría'': el que carga al niño en su seno & Nm 11:12 \\
\heb{אֹמֶנֶת} \emph{oménet} & ``aya'', nodriza: Noemí con el hijo de Rut & Rt 4:16 \\
\heb{אֹמְנוֹת} \emph{omnot} & ``quiciales'', los soportes de las puertas del templo & 2~R 18:16 \\
\heb{אֱמוּנָה} \emph{emuná} & ``firmeza'' de las manos de Moisés sostenidas por Aarón y Hur & Éx 17:12 \\
\heb{נֶאֱמָן} \emph{ne'emán} & ``lugar firme'' donde se clava una estaca & Is 22:23 \\
\heb{וְנֶאְמַן} \emph{vene'mán} & ``será afirmada'' la casa de David & 2~S 7:16 \\
\bottomrule
\end{tabular}\par
{\footnotesize\color{gris}\textit{Tabla 2.} Usos concretos de la raíz \emph{amán}. \textit{Nota.} Textos de la RVR1960 y la BHS.}
\end{center}

\begin{verso}{Éxodo 17:12}
Y las manos de Moisés se cansaban; por lo que tomaron una piedra, y la pusieron debajo de él, y se sentó sobre ella; y Aarón y Hur sostenían sus manos, el uno de un lado y el otro de otro; así hubo en sus manos firmeza hasta que se puso el sol.
\end{verso}

En todos estos casos la raíz habla de sostener y de ser sostenido: el brazo que carga a un niño, el poste que soporta una puerta, las manos que no se abaten. Verdadero es, entonces, lo que resiste el peso que se le pone encima.

\precision{el sentido de una palabra lo fija su uso en contexto, no su etimología; deducir el significado a partir de la raíz es incurrir en la ``falacia de la raíz'' (Carson, 2013). Aquí los usos concretos confirman la idea de firmeza que el material de la serie atribuye a \emph{emet}, pero en cada texto es el contexto el que decide si \emph{emet} significa ``verdad'', ``fidelidad'' o ``seguridad''.}

\subsection{Creer es apoyarse en lo firme}
Isaías aprovecha esa misma raíz en un juego de palabras célebre:

\hebverso{אִם לֹא תַאֲמִינוּ כִּי לֹא תֵאָמֵנוּ}
\begin{verso}{Isaías 7:9}
Si vosotros no creyereis, de cierto no permaneceréis.
\end{verso}

\begin{filologia}{Is 7:9}
\heb{תַאֲמִינוּ} \emph{ta'aminu} & \heb{אמן} & hifil: ``os apoyáis, creéis''. Creer es apoyar el peso en algo firme. \\
\heb{תֵאָמֵנוּ} \emph{te'amenu} & \heb{אמן} & nifal: ``seréis sostenidos, quedaréis firmes''. \\
\end{filologia}

``Si no os afirmáis, no seréis afirmados.'' Verdad, fe y fidelidad brotan de un mismo tronco: \emph{emet} es lo que sostiene; \emph{he'emín}, el acto de apoyarse en ello; \emph{emuná}, la firmeza tanto de quien sostiene como de quien se apoya. De ahí que en la Biblia la fe no se oponga a la verdad, sino que descanse en lo verdadero.

\precision{(textual) la Septuaginta lee aquí otra cosa: \grc{καὶ ἐὰν μὴ πιστεύσητε, οὐδὲ μὴ συνῆτε}, ``y si no creéis, tampoco entenderéis''. El juego de palabras solo existe en hebreo.}

\subsection{El espectro completo}
La correspondencia entre las lenguas no es de uno a uno. En la RVR1960 una sola palabra española, ``verdad'', traduce varios términos del original, y una sola raíz hebrea se vierte con varias palabras españolas. ``Verdad'' y sus formas (``verdadero'', ``verdaderamente'', ``veraz'') aparecen 432 veces en 406 versículos de esa versión: 188 en el Antiguo Testamento y 244 en el Nuevo.

\begin{center}\small\renewcommand{\arraystretch}{1.16}
\begin{tabular}{@{}L{3.6cm}L{1.1cm}L{5.4cm}L{4.2cm}@{}}
\toprule
\textbf{Término} & \textbf{Veces} & \textbf{Traducción en la RVR1960} & \textbf{Textos} \\ \midrule
\heb{אֱמֶת} \emph{emet} & 127 & ``verdad'' ($\approx$\,93), ``verdadero/a'' (9), ``fidelidad'' (4), ``seguridad'', ``rectamente'' & Gn 24:27; Jer 10:10; Sal 31:5 \\
\heb{אֱמוּנָה} \emph{emuná} & 49 & ``verdad'' (18), ``fidelidad'' (13), ``firmeza'', ``fe'' & Éx 17:12; Dt 32:4; Hab 2:4 \\
\heb{אָמֵן} \emph{amén} & 30 & ``amén'' (28); ``Dios de verdad'' (2) & Dt 27:15; Is 65:16 \\
\heb{אֹמֶן} \emph{ómen} & 1 & ``firmeza'' & Is 25:1 \\
\heb{אָמְנָם}, \heb{אֻמְנָם} \emph{omnam} & 14 & ``verdaderamente'', ``de veras'' & 1~R 8:27; Job 19:4 \\
\heb{אָכֵן} \emph{ajén} & --- & ``verdaderamente'' & Is 45:15 \\
\heb{נָכוֹן} \emph{najón} & --- & ``cosa cierta'' (junto a \emph{emet}) & Dt 13:14; 17:4 \\
\heb{קֹשֶׁט} \emph{qóshet} & 2 & ``verdad'' & Pr 22:21 \\
arameo \heb{קְשֹׁט} \emph{qeshot} & 2 & ``verdaderas'', ``ciertamente'' & Dn 2:47; 4:37 \\
arameo \heb{יַצִּיב} \emph{yatsiv} & 5 & ``verdadero'', ``la verdad'' & Dn 2:45; 7:16 \\ \midrule
\grc{ἀλήθεια} \emph{alétheia} & 109 & ``verdad'' (104) & Jn 14:6; Ro 1:25 \\
\grc{ἀληθής} \emph{alēthés} & 26 & ``verdadero'', ``veraz'' & Jn 8:26; Ro 3:4 \\
\grc{ἀληθινός} \emph{alēthinós} & 28 & ``verdadero'' (genuino, real) & Jn 15:1; 1~Ts 1:9 \\
\grc{ἀληθῶς} \emph{alēthós} & 18 & ``verdaderamente'' & 1~Ts 2:13 \\
\grc{ἀληθεύω} \emph{alētheúō} & 2 & ``decir la verdad'', ``seguir la verdad'' & Gá 4:16; Ef 4:15 \\
\grc{ἀψευδής} \emph{apseudés} & 1 & ``que no miente'' & Tit 1:2 \\
\grc{ἀμήν} \emph{amén} & 128 & ``de cierto'', ``amén'' & 2~Co 1:20; Ap 3:14 \\ \midrule
\heb{שֶׁקֶר} \emph{shéqer} & 113 & ``mentira, falsedad, engaño'' & Jer 10:14; 28:15 \\
\heb{כָּזָב} \emph{kazav} & 31 & ``mentira''; el verbo: ``mentir'' & Sal 116:11; Is 28:17 \\
\grc{ψεῦδος}, \grc{ψεύστης} & 10; 10 & ``mentira''; ``mentiroso'' & Ro 1:25; Jn 8:44 \\
\bottomrule
\end{tabular}\par
{\footnotesize\color{gris}\textit{Tabla 3.} El espectro de ``verdad''. \textit{Nota.} Conteo propio con la BHS y NA27 con números Strong, y la RVR1960 sin títulos de sección.}
\end{center}

\precision{los conteos son propios y conviene verificarlos en una concordancia (e-Sword). Para \emph{ajén} y \emph{najón} no se da cifra, porque en la mayoría de sus usos no significan ``verdad''. Tampoco se cuenta el giro ``a la verdad'', que la RVR1960 usa 57 veces: casi siempre traduce la partícula griega \grc{μέν} (``por cierto'') y no trata de la verdad como tema.}

\subsection{Una palabra, dos traducciones}
Puesto que \emph{emet} y \emph{emuná} reúnen verdad y fidelidad, las versiones se inclinan unas veces por un sentido y otras por el otro, como deja ver el compendio del material de la serie:

\begin{center}\small\renewcommand{\arraystretch}{1.2}
\begin{tabular}{@{}L{1.8cm}L{2.4cm}L{3.3cm}L{3.2cm}L{3.4cm}@{}}
\toprule
\textbf{Texto} & \textbf{Original} & \textbf{RVR1960} & \textbf{LBLA} & \textbf{NVI} \\ \midrule
Sal 119:30 & \heb{דֶּרֶךְ־אֱמוּנָה} & ``el camino de la verdad'' & ``el camino de la verdad'' & ``el camino de la fidelidad'' \\
Sal 119:160 & \heb{רֹאשׁ־דְּבָרְךָ אֱמֶת} & ``La suma de tu palabra es verdad'' & ``La suma de tu palabra es verdad'' & ``La suma de tus palabras es la verdad'' \\
Jer 5:1 & \heb{מְבַקֵּשׁ אֱמוּנָה} & ``que busque verdad'' & ``que busque la verdad'' & ``y busque la verdad'' \\
3~Jn 3 & \grc{τῇ ἀληθείᾳ} & ``tu verdad'' & ``tu verdad'' & ``tu fidelidad'' \\
\bottomrule
\end{tabular}\par
{\footnotesize\color{gris}\textit{Tabla 4.} Verdad y fidelidad en tres versiones. \textit{Nota.} NVI y LBLA según el compendio de la serie y la biblioteca bíblica.}
\end{center}

Las versiones no se contradicen: cada una destaca uno de los dos rostros de la misma raíz, porque en hebreo ser verdadero equivale a ser confiable. Sal 119:160 añade un matiz propio. \heb{רֹאשׁ} (\emph{rosh}) significa ``cabeza'' y admite dos lecturas: ``la suma'', es decir, el total, como cuando se cuentan cabezas, o ``el principio''. La RVR1960 traduce ``La suma de tu palabra es verdad''; la RV1909 y el material de la serie, ``El principio de tu palabra es verdad''. Ambas afirman que la palabra de Dios es verdad de punta a punta.

\subsection{La primera vez: el siervo de Abraham}
Cuando \emph{emet} aparece por primera vez en la Biblia, lo hace en pareja con \emph{jésed}. El siervo de Abraham, enviado a buscar esposa para Isaac, pide a Dios una señal (Gn 24:12--14), y la respuesta llega tal como la había pedido:

\hebverso{בָּרוּךְ יְהוָה אֱלֹהֵי אֲדֹנִי אַבְרָהָם אֲשֶׁר לֹא־עָזַב חַסְדּוֹ וַאֲמִתּוֹ מֵעִם אֲדֹנִי}
\begin{verso}{Génesis 24:27}
y dijo: Bendito sea \SENOR{}, Dios de mi amo Abraham, que no apartó de mi amo su misericordia y su verdad, guiándome \SENOR{} en el camino a casa de los hermanos de mi amo.
\end{verso}

Al relatar lo ocurrido, el siervo repite que Dios lo guió ``por camino de verdad'' (\heb{בְּדֶרֶךְ אֱמֶת}, Gn 24:48), esto es, por el camino correcto, el que llega a destino. Desde su primera aparición, pues, la verdad es la de un Dios que guía y cumple.

\subsection{El Dios del Amén}\label{sec:amen-verdad}
\hebverso{הַמִּתְבָּרֵךְ בָּאָרֶץ יִתְבָּרֵךְ בֵּאלֹהֵי אָמֵן \\ וְהַנִּשְׁבָּע בָּאָרֶץ יִשָּׁבַע בֵּאלֹהֵי אָמֵן}
\begin{verso}{Isaías 65:16}
El que se bendijere en la tierra, en el Dios de verdad se bendecirá; y el que jurare en la tierra, por el Dios de verdad jurará; porque las angustias primeras serán olvidadas, y serán cubiertas de mis ojos.
\end{verso}

\begin{filologia}{Is 65:16}
\heb{אֱלֹהֵי אָמֵן} \emph{Elohé amén} & \heb{אמן} & literalmente ``el Dios del Amén''. La RVR1960 traduce ``Dios de verdad''. \\
\grc{τὸν θεὸν τὸν ἀληθινόν} & --- & la Septuaginta traduce ``el Dios verdadero'', con el adjetivo que Pablo usará en 1~Ts 1:9. \\
\end{filologia}

La frase pertenece al anuncio de los cielos nuevos y la tierra nueva (Is 65:17), y en ella Dios pone su ``amén'' sobre su propia palabra: lo que promete, es. Así se prepara el título que el Apocalipsis dará a Cristo: ``He aquí el Amén, el testigo fiel y verdadero'' (Ap 3:14).

\subsection{Lo que permanece y lo que dura un momento}
\hebverso{שְׂפַת־אֱמֶת תִּכּוֹן לָעַד וְעַד־אַרְגִּיעָה לְשׁוֹן שָׁקֶר}
\begin{verso}{Proverbios 12:19}
El labio veraz permanecerá para siempre; Mas la lengua mentirosa sólo por un momento.
\end{verso}

\begin{filologia}{Pr 12:19}
\heb{שְׂפַת־אֱמֶת} \emph{sefat emet} & \heb{אמן} & ``labio de verdad''. \\
\heb{תִּכּוֹן} \emph{tikkón} & \heb{כון} & ``estará firme, establecido''. \\
\heb{עַד־אַרְגִּיעָה} \emph{ad argi'á} & \heb{רגע} & ``hasta que yo parpadee'': un instante. \\
\end{filologia}

El proverbio mide la verdad por su duración: lo verdadero queda en pie, mientras la mentira dura lo que un parpadeo. Es la plomada aplicada al tiempo.

\subsection{Del hebreo a Pablo}
La Septuaginta traduce \emph{emet} casi siempre con \grc{ἀλήθεια}, vierte ``el Dios del Amén'' de Is 65:16 como \grc{τὸν θεὸν τὸν ἀληθινόν} y conserva \emph{amén} sin traducir en las fórmulas litúrgicas. Ese es el vocabulario que Pablo hereda cuando llama a Dios \grc{ἀληθής}, ``veraz'' (Ro 3:4), y \grc{ἀληθινός}, ``verdadero'' (1~Ts 1:9), y cuando afirma que en Cristo está ``el Amén'' de todas las promesas (2~Co 1:20).

\begin{center}
\begin{tikzpicture}[font=\small,node distance=0.9cm]
 \node[draw=azul,line width=0.8pt,fill=papiro,rounded corners=3pt,text width=3.6cm,align=center,inner sep=6pt] (h) {\textbf{Hebreo}\\ \heb{אֱמֶת} \emph{emet}\\ \heb{אָמֵן} \emph{amén}\\ \heb{אֱמוּנָה} \emph{emuná}};
 \node[draw=azul,line width=0.8pt,fill=papiro,rounded corners=3pt,text width=3.6cm,align=center,inner sep=6pt,right=of h] (l) {\textbf{Septuaginta}\\ \grc{ἀλήθεια}\\ \grc{ἀληθινός} (Is 65:16)\\ \grc{πίστις}, \grc{ἀλήθεια}};
 \node[draw=dorado,line width=1.6pt,fill=azul,text=papiro,rounded corners=3pt,text width=3.6cm,align=center,inner sep=6pt,right=of l] (p) {\textbf{Pablo}\\ \grc{ὁ θεὸς ἀληθής} (Ro 3:4)\\ \grc{θεῷ\ldots{} ἀληθινῷ} (1 Ts 1:9)\\ \grc{τὸ ἀμήν} (2 Co 1:20)};
 \draw[-{Latex[length=3mm]},dorado,line width=1.4pt] (h)--(l);
 \draw[-{Latex[length=3mm]},dorado,line width=1.4pt] (l)--(p);
\end{tikzpicture}\par
{\footnotesize\color{gris}\textit{Figura 1.} Del hebreo a Pablo. \textit{Nota.} Elaboración propia.}
\end{center}

En griego clásico, \grc{ἀλήθεια} designa lo real por oposición a la apariencia. En la Septuaginta y en Pablo, sin embargo, la palabra absorbe además el sentido hebreo de firmeza y fidelidad, hasta el punto de que en Ro 3:3--7 ``la fidelidad de Dios'' y ``la verdad de Dios'' nombran una misma realidad.

\precision{suele explicarse \grc{ἀλήθεια} como \emph{a-lēthē}, ``lo no oculto''. Es una etimología discutida que, en cualquier caso, no decide el sentido bíblico de la palabra, pues este depende del uso (Carson, 2013).}

%===================================================================
\section{El Nombre que revela la verdad}\label{sec:nombre-verdad}

En la Escritura el nombre de Dios no es una etiqueta, sino la revelación de lo que él es. Cuando Moisés pidió ver la gloria, Dios respondió que haría pasar su bien delante de él y proclamaría su Nombre (Éx 33:18--19); y al proclamarlo se declaró ``grande en misericordia y verdad'' (\heb{וְרַב־חֶסֶד וֶאֱמֶת}, \emph{verav jésed ve'emet}, Éx 34:6). Desde el Sinaí, por tanto, la verdad forma parte del Nombre. Dos de los títulos que la Escritura forma con la raíz \emph{amán} ya han aparecido en este capítulo: \emph{Elohim emet}, ``el Dios verdadero'' del texto ancla (Jer 10:10), y \emph{Elohé amén}, ``el Dios del Amén'' (Is 65:16, véase la subsección~\ref{sec:amen-verdad}). Quedan otros dos, nacidos ambos en horas de crisis: uno en labios de Moisés, a las puertas de la muerte; el otro en labios de David, acosado por sus enemigos.

\subsection{\emph{El emuná}: el Nombre que Moisés proclamó antes de morir}
El cántico de Deuteronomio 32 es la última palabra pública de Moisés. Dios le había anunciado: ``He aquí se ha acercado el día de tu muerte'' (Dt 31:14), y le había mandado escribir un cántico que quedara ``por testigo contra los hijos de Israel'' (31:19), porque el pueblo, una vez saciado en la tierra, se volvería a otros dioses. Moisés no se hacía ilusiones: ``yo sé que después de mi muerte, ciertamente os corromperéis'' (31:29). El cántico nace, pues, de la certeza de la infidelidad del pueblo, y su primer gesto es proclamar un Nombre:

\hebverso{כִּי שֵׁם יְהוָה אֶקְרָא הָבוּ גֹדֶל לֵאלֹהֵינוּ \\ הַצּוּר תָּמִים פָּעֳלוֹ כִּי כָל־דְּרָכָיו מִשְׁפָּט אֵל אֱמוּנָה וְאֵין עָוֶל צַדִּיק וְיָשָׁר הוּא}
\begin{verso}{Deuteronomio 32:3–4}
Porque el nombre de \SENOR{} proclamaré. Engrandeced a nuestro Dios. El es la Roca, cuya obra es perfecta, Porque todos sus caminos son rectitud; Dios de verdad, y sin ninguna iniquidad en él; Es justo y recto.
\end{verso}

\begin{filologia}{Dt 32:3--4}
\heb{שֵׁם יְהוָה אֶקְרָא} \emph{shem Adonay eqrá} & \heb{קרא} & ``el nombre del \SENOR{} proclamaré'': el mismo verbo de Éx 34:5--6, donde Dios ``proclamó'' (\heb{וַיִּקְרָא}) su propio Nombre. Moisés repite al final de su vida lo que oyó en el Sinaí. \\
\heb{הַצּוּר} \emph{hatsur} & \heb{צור} & ``la Roca'': título de Dios a lo largo del cántico (Dt 32:4, 15, 18, 30, 31). Lo que no se mueve. \\
\heb{אֵל אֱמוּנָה} \emph{El emuná} & \heb{אמן} & ``Dios de \emph{emuná}'': de firmeza, de fidelidad. La RVR1960 traduce ``Dios de verdad''; la LXX, \grc{θεὸς πιστός}, ``Dios fiel''. \\
\grc{θεός, ἀληθινὰ τὰ ἔργα αὐτοῦ} & --- & la LXX vierte ``cuya obra es perfecta'' como ``Dios, verdaderas son sus obras'': pone la verdad también al comienzo del verso. \\
\end{filologia}

\emph{El emuná} se comprende por contraste. Pocos versos después, el mismo cántico retrata al pueblo como ``una generación perversa, hijos infieles'' (Dt 32:20), y en hebreo ``infieles'' es \heb{בָּנִים לֹא־אֵמֻן בָּם}, ``hijos en quienes no hay \emph{emun}'': la misma raíz, ahora negada. La Septuaginta hace el contraste todavía más visible, pues llama a Dios \grc{θεὸς πιστός} (32:4) y a los hijos \grc{υἱοὶ οἷς οὐκ ἔστιν πίστις ἐν αὐτοῖς}, ``hijos en quienes no hay fidelidad'' (32:20). La verdad de Dios no se edifica, entonces, sobre la del pueblo, sino a pesar de ella. Moisés no puede legar a Israel un pueblo firme, pero sí un Nombre firme; por eso el título va unido a la Roca, porque lo que Israel necesitará cuando todo se le derrumbe no es su propia constancia, sino un Dios que no se mueve. La subsección~\ref{sec:recto-verdad} expone esta misma imagen leída como plomada; aquí interesa que se trata de un nombre y que lo proclamó un hombre que iba a morir sin entrar en la tierra (Dt 32:48--52).

\subsection{\emph{El emet}: el Nombre en que David depositó su vida}
El Salmo 31 lleva el título ``de David'' (\heb{לְדָוִד}), y su autor habla como un hombre cercado: ``oigo la calumnia de muchos; el miedo me asalta por todas partes, mientras consultan juntos contra mí e idean quitarme la vida'' (Sal 31:13 [31:14]). En medio de esa angustia se entrega:

\hebverso{בְּיָדְךָ אַפְקִיד רוּחִי \\ פָּדִיתָה אוֹתִי יְהוָה אֵל אֱמֶת}
\begin{verso}{Salmo 31:5 [31:6]}
En tu mano encomiendo mi espíritu; Tú me has redimido, oh \SENOR{}, Dios de verdad.
\end{verso}

\begin{filologia}{Sal 31:5 (31:6 en hebreo)}
\heb{אַפְקִיד} \emph{afqid} & \heb{פקד} & ``encomiendo'', ``deposito'': entregar algo en custodia. De la misma raíz es \heb{פִּקָּדוֹן} (\emph{piqqadón}), ``el depósito que se le encomendó'' (Lv 6:4 [5:23]). \\
\heb{פָּדִיתָה} \emph{paditá} & \heb{פדה} & ``me has redimido'': rescatar pagando un precio. El perfecto mira hacia atrás: la redención pasada funda la entrega presente. \\
\heb{אֵל אֱמֶת} \emph{El emet} & \heb{אמן} & ``Dios de verdad'': el Dios firme, en cuyas manos lo depositado está seguro. LXX \grc{κύριε ὁ θεὸς τῆς ἀληθείας} (Sal 30:6 LXX). \\
\end{filologia}

El verbo \emph{afqid} pertenece al lenguaje de los depósitos: la Ley regulaba qué debía hacerse cuando alguien recibía en custodia un bien ajeno y no lo devolvía (Lv 6:2--4 [5:21--23]). David deposita lo más precioso que posee, su propio aliento, en manos de un depositario que no puede fallar, y el nombre \emph{El emet} es la garantía de ese depósito, porque solo se encomienda algo a quien es firme. El verso siguiente traza el contraste con la misma lógica de Jeremías 10: ``Aborrezco a los que esperan en vanidades ilusorias; mas yo en \SENOR{} he esperado'' (Sal 31:6 [31:7]). Las ``vanidades'' son \heb{הַבְלֵי־שָׁוְא}, \emph{havlé shav}, con la misma palabra \emph{hével} con que Jeremías designa a los ídolos. De un lado queda el Dios de verdad; del otro, el soplo.

\precision{el título no indica la ocasión del salmo. Algunos lectores la sitúan en la persecución de Saúl, porque el ``en mi premura'' de Sal 31:22 [31:23] (\heb{בְחָפְזִי}) emplea la raíz de 1~S 23:26, ``se daba prisa David para escapar de Saúl'' (\heb{נֶחְפָּז}). La conexión verbal es sugerente, pero no permite una identificación segura.}

\subsection{El Nombre en los labios de Jesús}
Siglos después, Jesús pronunció ese mismo verso en la cruz:

\begin{verso}{Lucas 23:46}
Entonces Jesús, clamando a gran voz, dijo: Padre, en tus manos encomiendo mi espíritu. Y habiendo dicho esto, expiró.
\end{verso}

Lucas cita la primera línea del salmo casi con las palabras de la Septuaginta, y así \grc{εἰς χεῖράς σου παραθήσομαι τὸ πνεῦμά μου}, ``en tus manos encomendaré mi espíritu'' (Sal 30:6 LXX), se convierte en \grc{πάτερ, εἰς χεῖράς σου παρατίθεμαι τὸ πνεῦμά μου}. Los cambios son dos: el verbo pasa al presente, porque la entrega ocurre en ese mismo instante, y se antepone ``Padre''. Lo que David hizo en peligro de muerte, Jesús lo hace en la muerte misma, ante el mismo depositario: el Dios que no falla con lo que se le confía y que al tercer día lo levantó de los muertos.

\precision{Lucas cita solo la primera línea del verso. Que Jesús se entregara al \emph{El emet} de la segunda línea es una lectura propia, que se apoya en que la cita toma el comienzo del versículo y en que el oyente judío conocía su final.}

La figura se prolonga en el resto del Nuevo Testamento. Pedro exhorta a los que sufren a que ``encomienden sus almas al fiel Creador'' (1~P 4:19), con el verbo de Lucas (\grc{πιστῷ κτίστῃ παρατιθέσθωσαν}); y Pablo, preso y cercano a su propia muerte, escribe: ``yo sé a quién he creído, y estoy seguro que es poderoso para guardar mi depósito para aquel día'' (2~Ti 1:12), donde ``depósito'' traduce \grc{παραθήκη}, de la misma familia. Cristo depositó su vida en el \emph{El emet}, y los que están en él depositan la suya en ese mismo Dios.

\subsection{El Verdadero y el Amén, con Pablo}
Lo que Moisés y David nombraron, el Nuevo Testamento lo halla en una Persona. El Resucitado se presenta a la iglesia de Filadelfia como ``el Santo, el Verdadero'' (\grc{ὁ ἅγιος, ὁ ἀληθινός}, Ap 3:7), y a la de Laodicea como ``el Amén, el testigo fiel y verdadero'' (\grc{ὁ ἀμήν, ὁ μάρτυς ὁ πιστὸς καὶ ἀληθινός}, Ap 3:14). Este último título reúne las tres palabras de la raíz: el \emph{amén} de Is 65:16, la fidelidad (\emph{emuná}, \grc{πιστός}) de Dt 32:4 y la verdad (\emph{emet}, \grc{ἀληθινός}) de Sal 31:5. Jesús ya lo había dicho de sí mismo: ``Yo soy el camino, y la verdad, y la vida'' (Jn 14:6).

Pablo lee los nombres del mismo modo. Ante la objeción de que la infidelidad de Israel anularía la fidelidad de Dios, responde: ``sea Dios veraz, y todo hombre mentiroso'' (Ro 3:4), que reproduce exactamente el contraste de Dt 32:4 y 32:20; y afirma que todas las promesas de Dios ``son en él Sí, y en él Amén'' (2~Co 1:20). Ambos textos se exponen en la sección dedicada a la verdad en Cristo (sección~\ref{sec:cristo-verdad}). A Tito, por su parte, le recuerda que la esperanza de la vida eterna la prometió ``Dios, que no miente'' (\grc{ὁ ἀψευδὴς θεός}, Tit 1:2), que es la forma negativa del \emph{El emet}.

\begin{center}\small\renewcommand{\arraystretch}{1.22}
\begin{tabular}{@{}L{3.3cm}L{3.7cm}L{2.4cm}L{5.0cm}@{}}
\toprule
\textbf{Nombre o título} & \textbf{Hebreo / griego} & \textbf{Texto} & \textbf{Qué revela} \\ \midrule
Grande en misericordia y verdad & \heb{וְרַב־חֶסֶד וֶאֱמֶת} & Éx 34:6 & la verdad está dentro del Nombre proclamado en el Sinaí \\
El Dios verdadero & \heb{אֱלֹהִים אֱמֶת} & Jer 10:10 & el Dios vivo frente al ídolo clavado \\
Dios de verdad (fidelidad) & \heb{אֵל אֱמוּנָה}; LXX \grc{θεὸς πιστός} & Dt 32:4 & la Roca firme frente a hijos ``sin \emph{emun}'' \\
Dios de verdad & \heb{אֵל אֱמֶת}; LXX \grc{ὁ θεὸς τῆς ἀληθείας} & Sal 31:5 [31:6] & en sus manos se deposita la vida \\
El Dios del Amén & \heb{אֱלֹהֵי אָמֵן}; LXX \grc{τὸν θεὸν τὸν ἀληθινόν} & Is 65:16 & dice ``amén'' a su propia palabra \\
Dios que no miente & \grc{ὁ ἀψευδὴς θεός} & Tit 1:2 & su promesa de vida eterna es segura \\
El Santo, el Verdadero & \grc{ὁ ἅγιος, ὁ ἀληθινός} & Ap 3:7 & Cristo, el Dios genuino \\
El Amén, el testigo fiel y verdadero & \grc{ὁ ἀμήν, ὁ μάρτυς ὁ πιστὸς καὶ ἀληθινός} & Ap 3:14 & en Cristo se reúnen \emph{amén}, \emph{emuná} y \emph{emet} \\
\bottomrule
\end{tabular}\par
{\footnotesize\color{gris}\textit{Tabla 5.} Los nombres que revelan la verdad. \textit{Nota.} Elaboración propia con la BHS, la LXX (Rahlfs-Hanhart) y NA28; numeración española [hebrea].}
\end{center}

\precision{la tradición rabínica condensó este principio al llamar a la verdad ``el sello del Santo, bendito sea'' (\emph{Shabat} 55a, dicho de Rabí Janiná citado por Resh Laquish). La sección~\ref{sec:judaismo-verdad} expone el texto y su explicación a partir de Jer 10:10.}

%===================================================================
\section{Definición}

\subsection{La verdad como correspondencia}
El material de la serie ofrece una definición directa, breve y concreta: ``La verdad es lo que corresponde a la realidad''. Una idea, un sentimiento o una estructura que no corresponda a la realidad tal como Dios la ve, la concibió y la diseñó carece de firmeza y, por eso mismo, no es verdad.

\precision{entender la verdad como correspondencia con la realidad es la definición clásica en filosofía. Lo propio de la mirada bíblica está en la medida: la realidad no es lo que percibe el observador, sino lo que Dios concibió, creó y sostiene. De ahí que la verdad humana sea siempre derivada, y la de Dios, la fuente.}

\begin{tcolorbox}[colback=papiro,colframe=azul,boxrule=1pt,arc=3pt,fonttitle=\bfseries,title=Definición]
La verdad de Dios es la perfección eterna por la cual Dios es real y firme, dice lo que es y cumple lo que dice. Su palabra corresponde exactamente a la realidad que él mismo creó y sostiene; por eso él es la medida de toda verdad (Jer 10:10; Nm 23:19; Jn 14:6).
\end{tcolorbox}

\begin{center}\small\renewcommand{\arraystretch}{1.2}
\begin{tabular}{@{}L{3.4cm}L{6.6cm}L{4.4cm}@{}}
\toprule
\textbf{Rasgo} & \textbf{Qué afirma} & \textbf{Textos} \\ \midrule
Real & es el Dios vivo, no una imagen & Jer 10:10; 1~Ts 1:9 \\
Firme (\emph{emet}) & lo que dice se sostiene & Sal 119:160; Pr 12:19 \\
Veraz & no miente ni puede mentir & Nm 23:19; Tit 1:2; He 6:18 \\
Fiel (\emph{emuná}) & cumple sus promesas & Mi 7:20; Ro 15:8 \\
Una & no dice ``sí'' y ``no'' a la vez & 2~Co 1:18--19 \\
Encarnada & la verdad tiene rostro & Jn 1:14; 14:6 \\
\bottomrule
\end{tabular}\end{center}

\textbf{Clasificación.} La verdad se cuenta entre los atributos de la \emph{relación con el ser humano}, pues responde a la pregunta ``¿cómo se revela Dios?''. Como todas sus perfecciones, es eterna (``su verdad por todas las generaciones'', Sal 100:5). En cuanto fuente, es \emph{incomunicable}: solo Dios es la verdad, y ningún otro tiene verdad como la suya. En cuanto veracidad, es \emph{comunicable}: ``hablad verdad cada uno con su prójimo'' (Ef 4:25).

\subsection{La verdad es una}
La verdad es singular, no plural. El cristianismo no propone ``múltiples verdades'', sino una verdad escrita con ``V'' mayúscula, que se mantiene unida y ofrece una visión del mundo coherente consigo misma. Pablo lo afirma de la palabra de Dios y de su propia predicación: ``como Dios es fiel, nuestra palabra a vosotros no es Sí y No'' (2~Co 1:18). Y Jesús lo pone en su oración: ``Santifícalos en tu verdad; tu palabra es verdad'' (Jn 17:17).

\subsection{Dios es luz: una verdad sin doblez}
Juan condensa en una sola frase el mensaje que recibió de Cristo:

\begin{verso}{1 Juan 1:5–7}
Este es el mensaje que hemos oído de él, y os anunciamos: Dios es luz, y no hay ningunas tinieblas en él. Si decimos que tenemos comunión con él, y andamos en tinieblas, mentimos, y no practicamos la verdad; pero si andamos en luz, como él está en luz, tenemos comunión unos con otros, y la sangre de Jesucristo su Hijo nos limpia de todo pecado.
\end{verso}

\begin{filologia}{1 Jn 1:5--6}
\grc{ὁ θεὸς φῶς ἐστιν} & \grc{φῶς} & ``Dios es luz'': no dice que Dios tenga luz, sino que \emph{es} luz, como Jer 10:10 dice que \emph{es} verdad. \\
\grc{σκοτία ἐν αὐτῷ οὐκ ἔστιν οὐδεμία} & \grc{σκοτία} & ``tiniebla en él no hay, ninguna'': la doble negación excluye hasta la mínima sombra. \\
\grc{ψευδόμεθα} & \grc{ψεύδομαι} & ``mentimos'': decir comunión con la luz y andar en tinieblas. \\
\grc{οὐ ποιοῦμεν τὴν ἀλήθειαν} & \grc{ἀλήθεια} & ``no hacemos la verdad'': la verdad se practica, no solo se dice (cf. Jn 3:21). \\
\end{filologia}

Que Dios sea luz sin tiniebla alguna funda una verdad \emph{absoluta}: no hay en él una zona oscura en la que su palabra deje de valer. Santiago añade que en el ``Padre de las luces'' no hay ``mudanza, ni sombra de variación'' (Stg 1:17): \grc{παρ᾽ ᾧ οὐκ ἔνι παραλλαγὴ ἢ τροπῆς ἀποσκίασμα}. El sol proyecta sombras que cambian a medida que gira; Dios no, y su verdad no varía con la hora ni con el ángulo desde el que se la mire.

La misma luz funda, además, una verdad \emph{sin doblez}, y el salmista retrata la mentira precisamente como un corazón dividido:

\hebverso{שָׁוְא יְדַבְּרוּ אִישׁ אֶת־רֵעֵהוּ שְׂפַת חֲלָקוֹת בְּלֵב וָלֵב יְדַבֵּרוּ}
\begin{verso}{Salmo 12:2 [12:3]}
Habla mentira cada uno con su prójimo; Hablan con labios lisonjeros, y con doblez de corazón.
\end{verso}

``Con doblez de corazón'' traduce \heb{בְּלֵב וָלֵב} (\emph{belev valev}), literalmente ``con un corazón y un corazón'', y Santiago llama a quien vive así \grc{δίψυχος}, ``de doble alma'' (Stg 1:8). Dios, en cambio, no es ``Sí y No'' (2~Co 1:17--18), y Jesús pide a los suyos esa misma sencillez: ``sea vuestro hablar: Sí, sí; no, no'' (Mt 5:37). Pablo extiende la exigencia a la iglesia: ``ahora sois luz en el Señor; andad como hijos de luz (porque el fruto del Espíritu es en toda bondad, justicia y verdad)'' (Ef 5:8--9). Los Salmos ya habían unido ambas palabras: ``Envía tu luz y tu verdad; éstas me guiarán'' (Sal 43:3).

\precision{(textual) en Ef 5:9 el texto crítico (NA28) lee \grc{ὁ καρπὸς τοῦ φωτός}, ``el fruto de la luz''; la RVR1960 sigue el Texto Recibido, ``el fruto del Espíritu''. La lectura de NA28 une directamente la luz con la verdad.}

\subsection{Antídoto al relativismo}
Una verdad absoluta y eterna como la de Dios se convierte en antídoto contra el relativismo moderno, porque excluye que la verdad dependa de la cultura, del individuo o del pragmatismo (``lo que funciona''). Es tan absoluta como la muerte o como la gravedad, y no cambia por lo que opine la sociedad: lo que Dios calificó de profano en el primer siglo lo sigue siendo hoy.

\begin{verso}{Salmo 119:89}
Para siempre, oh \SENOR{}, Permanece tu palabra en los cielos.
\end{verso}
\begin{verso}{Salmo 119:160}
La suma de tu palabra es verdad, Y eterno es todo juicio de tu justicia.
\end{verso}

Miguel Núñez recuerda que Francis Schaeffer, ante la idea de que cada uno tiene su propia verdad, acuñó la expresión ``la verdadera verdad'', ``dando a entender que la verdad no es nuestra, sino que es únicamente de Dios'' (Núñez, 2020, cap. 1).

%===================================================================
\section{La plomada: la verdad como medida}\label{sec:plomada}

\subsection{Cómo funciona una plomada}
La imagen central del material de la serie es la plomada del albañil: un peso que cuelga de un cordel y que la gravedad sitúa siempre en la vertical exacta, opine lo que opine quien construye. No es la plomada la que se ajusta al muro, sino el muro el que se ajusta a ella, y el albañil que la ignora levanta una pared que puede parecer firme, pero que tarde o temprano se derrumba.

La verdad de Dios funciona del mismo modo. Es la plomada universal, que mide la verticalidad exacta de lo que existe, de lo que creemos y de cómo nos conducimos. Lo que no se alinea con ella es, por definición, mentira, y está destinado al colapso; y, al igual que la gravedad, no depende de lo que la sociedad piense de ella.

\begin{center}
\begin{tikzpicture}[scale=0.95]
 \fill[gris!25] (-6.2,-0.15) rectangle (6.4,0);
 % muro a plomo
 \foreach \y in {0,0.45,...,3.2}{\draw[azul,fill=papiro] (-6.0,\y) rectangle (-4.9,\y+0.45);}
 \node[font=\scriptsize\bfseries,text=azul,align=center] at (-5.45,-0.45) {a plomo: \emph{emet}};
 % plomada
 \fill[azuloscuro] (-3.4,4.1) rectangle (-1.9,4.3);
 \draw[azuloscuro,line width=1pt] (-2.65,4.1) -- (-2.65,0.75);
 \fill[dorado] (-2.87,0.75) -- (-2.43,0.75) -- (-2.65,0.2) -- cycle;
 \node[font=\scriptsize\bfseries,text=azuloscuro] at (-2.65,4.6) {la plomada};
 \node[font=\scriptsize,text=azul,align=center,text width=5cm] at (-2.65,-1.15) {la verdad de Dios: ``La suma de tu palabra es verdad'' (Sal 119:160)};
 % muro torcido
 \foreach \i in {0,...,7}{\pgfmathsetmacro{\dx}{0.11*\i}\draw[gris,fill=gris!15] (-1.0+\dx,0.45*\i) rectangle (0.1+\dx,0.45*\i+0.45);}
 \draw[gris,dashed] (-1.0,0) -- (-1.0,3.8);
 \node[font=\scriptsize\bfseries,text=gris,align=center] at (0.1,-0.45) {fuera de plomo: \emph{shéqer}};
 % mundo en deriva
 \draw[gris,line width=0.8pt] (2.0,3.8) -- (2.3,3.5) -- (1.9,3.2) -- (2.35,2.9) -- (2.0,2.6) -- (2.25,2.2) -- (1.8,1.9) -- (2.3,1.5) -- (1.95,1.2) -- (2.2,0.9) -- (1.9,0.55) -- (2.3,0.2);
 \node[anchor=west,align=left,text=gris,font=\scriptsize,text width=3.6cm] at (2.6,2.0) {\textbf{el mundo, torcido y en deriva}\\ ``generación torcida y perversa'' (Dt 32:5); ``llevados por doquiera de todo viento'' (Ef 4:14)};
\end{tikzpicture}\par\medskip
{\footnotesize\color{gris}\textit{Figura 2.} La plomada: la verticalidad de la verdad. \textit{Nota.} Elaboración propia a partir del material de la serie.}
\end{center}

\subsection{La plomada en la Escritura}
La plomada no es una ilustración moderna impuesta al texto; los profetas ya la usaban.

\begin{verso}{Amós 7:7–8}
Me enseñó así: He aquí el Señor estaba sobre un muro hecho a plomo, y en su mano una plomada de albañil. \SENOR{} entonces me dijo: ¿Qué ves, Amós? Y dije: Una plomada de albañil. Y el Señor dijo: He aquí, yo pongo plomada de albañil en medio de mi pueblo Israel; no lo toleraré más.
\end{verso}

\hebverso{הִנְנִי יִסַּד בְּצִיּוֹן אָבֶן אֶבֶן בֹּחַן פִּנַּת יִקְרַת מוּסָד מוּסָּד הַמַּאֲמִין לֹא יָחִישׁ \\ וְשַׂמְתִּי מִשְׁפָּט לְקָו וּצְדָקָה לְמִשְׁקָלֶת וְיָעָה בָרָד מַחְסֵה כָזָב}
\begin{verso}{Isaías 28:16–17}
por tanto, \SENOR{} el Señor dice así: He aquí que yo he puesto en Sion por fundamento una piedra, piedra probada, angular, preciosa, de cimiento estable; el que creyere, no se apresure. Y ajustaré el juicio a cordel, y a nivel la justicia; y granizo barrerá el refugio de la mentira, y aguas arrollarán el escondrijo.
\end{verso}

\begin{filologia}{Am 7:7--8 e Is 28:16--17}
\heb{אֲנָךְ} \emph{anaj} & --- & ``plomada'' (sus cuatro apariciones están en Am 7:7--8). \\
\heb{מִשְׁקֶלֶת} \emph{mishquélet} & \heb{שׁקל} & ``nivel, plomada'': el peso que cuelga del cordel (Is 28:17; 2~R 21:13). \\
\heb{קַו} \emph{qav} & \heb{קוה} & ``cordel'' de medir. \\
\heb{מוּסָד מוּסָּד} \emph{musad mussad} & \heb{יסד} & ``cimiento bien cimentado'' (RVR1960: ``de cimiento estable''). \\
\heb{הַמַּאֲמִין} \emph{hama'amín} & \heb{אמן} & ``el que creyere'': el que se apoya en la piedra firme. \\
\heb{מַחְסֵה כָזָב} \emph{majsé jazav} & \heb{כזב} & ``refugio de la mentira''. \\
\end{filologia}

Dios mide a su pueblo con su propia rectitud, y lo que no está a plomo se viene abajo: el granizo barre el ``refugio de la mentira''. El mismo oráculo presenta la piedra de cimiento junto con la raíz \emph{amán}, de modo que el que cree (\emph{hama'amín}) se apoya en lo que Dios dejó firme. Pablo leerá en esa piedra a Cristo (sección~\ref{sec:piedra}).

\precision{el sentido de \emph{anaj} es discutido: HALOT propone ``estaño'' (acadio \emph{anāku}) en lugar de ``plomada''. \emph{Mishquélet}, en cambio, designa con seguridad una plomada o un nivel. La imagen es, por tanto, firme en Is 28:17 y 2~R 21:13, y probable en Amós.}


\subsection{Recto y torcido}\label{sec:recto-verdad}
Moisés ya había cantado esta imagen en su último cántico, donde Dios es la vertical y el pueblo, lo torcido:

\hebverso{הַצּוּר תָּמִים פָּעֳלוֹ כִּי כָל־דְּרָכָיו מִשְׁפָּט אֵל אֱמוּנָה וְאֵין עָוֶל צַדִּיק וְיָשָׁר הוּא \\ שִׁחֵת לוֹ לֹא בָּנָיו מוּמָם דּוֹר עִקֵּשׁ וּפְתַלְתֹּל}
\begin{verso}{Deuteronomio 32:4–5}
El es la Roca, cuya obra es perfecta, Porque todos sus caminos son rectitud; Dios de verdad, y sin ninguna iniquidad en él; Es justo y recto. La corrupción no es suya; de sus hijos es la mancha, Generación torcida y perversa.
\end{verso}

\begin{filologia}{Dt 32:4--5}
\heb{אֵל אֱמוּנָה} \emph{El emuná} & \heb{אמן} & ``Dios de verdad'': Dios de firmeza. \\
\heb{יָשָׁר} \emph{yashar} & \heb{ישׁר} & ``recto'': derecho, a plomo. \\
\heb{עִקֵּשׁ} \emph{iqqesh} & \heb{עקשׁ} & ``torcido'', desviado de la línea. \\
\heb{פְּתַלְתֹּל} \emph{petaltol} & \heb{פתל} & ``retorcido'', enroscado. \\
LXX & --- & \grc{γενεὰ σκολιὰ καὶ διεστραμμένη}, ``generación torcida y pervertida''. \\
\end{filologia}

El Predicador resume toda la historia humana con esa misma palabra: ``Dios hizo al hombre recto (\heb{יָשָׁר}), pero ellos buscaron muchas perversiones'' (Ec 7:29). Pablo, por su parte, recurre al cántico de Moisés para describir la misión de la iglesia: ``para que seáis irreprensibles y sencillos, hijos de Dios sin mancha en medio de una generación maligna y perversa (\grc{γενεᾶς σκολιᾶς καὶ διεστραμμένης}), en medio de la cual resplandecéis como luminares en el mundo'' (Flp 2:15). Quienes han sido enderezados por la verdad se convierten, en medio de lo torcido, en la referencia visible.

\subsection{Todos se desviaron}
Medida con la plomada de Dios, la humanidad entera aparece fuera de la vertical:

\hebverso{הַכֹּל סָר יַחְדָּו נֶאֱלָחוּ אֵין עֹשֵׂה־טוֹב אֵין גַּם־אֶחָד}
\begin{verso}{Salmo 14:3}
Todos se desviaron, a una se han corrompido; No hay quien haga lo bueno, no hay ni siquiera uno.
\end{verso}

\heb{הַכֹּל סָר} significa ``todos se apartaron'', y el verbo \emph{sur} describe al que se sale del camino. Pablo incorpora el salmo a su diagnóstico de la humanidad: ``Todos se desviaron, a una se hicieron inútiles'' (Ro 3:12). Su verbo, \grc{ἐξέκλιναν}, quiere decir literalmente ``se inclinaron hacia afuera'': nadie quedó a plomo.

\precision{el material de la serie cita este salmo con otra versión (``se han vuelto inútiles'' corresponde a la forma de Ro 3:12); aquí se usa la RVR1960.}

\subsection{El muro que se cae de repente}
Isaías describe la suerte de un muro que ha perdido la vertical:

\hebverso{לָכֵן יִהְיֶה לָכֶם הֶעָוֹן הַזֶּה כְּפֶרֶץ נֹפֵל נִבְעֶה בְּחוֹמָה נִשְׂגָּבָה \\ אֲשֶׁר־פִּתְאֹם לְפֶתַע יָבוֹא שִׁבְרָהּ}
\begin{verso}{Isaías 30:12–13}
Por tanto, el Santo de Israel dice así: Porque desechasteis esta palabra, y confiasteis en violencia y en iniquidad, y en ello os habéis apoyado; por tanto, os será este pecado como grieta que amenaza ruina, extendiéndose en una pared elevada, cuya caída viene súbita y repentinamente.
\end{verso}

\begin{filologia}{Is 30:12--13}
\heb{וַתִּשָּׁעֲנוּ} \emph{vatisha'anu} & \heb{שׁען} & ``y os habéis apoyado'': lo contrario de \emph{he'emín}, apoyarse en lo firme (Is 7:9). \\
\heb{כְּפֶרֶץ נֹפֵל} \emph{keférets nofel} & \heb{פרץ} & ``como brecha que cae''. \\
\heb{נִבְעֶה} \emph{niv'é} & \heb{בעה} & ``abombado, que sobresale'': el muro que ya se salió de la vertical. \\
\heb{פִּתְאֹם לְפֶתַע} \emph{pitom lefeta} & --- & ``de repente, en un instante''. \\
\end{filologia}

El pecado es un muro fuera de plomo: parece firme hasta el día en que se desploma, y entonces cae de golpe. Es la lógica de Pr 12:19, según la cual la mentira dura ``un parpadeo''.

\subsection{Tres pasos de la plomada}
\begin{enumerate}
\item \textbf{Dios cuelga la plomada.} Es él, y no nosotros, quien fija la medida: ``He aquí, yo pongo plomada de albañil en medio de mi pueblo'' (Am 7:8); ``La suma de tu palabra es verdad'' (Sal 119:160).
\item \textbf{Se mide el muro.} Creencias, sentimientos y conductas se examinan a su luz: ``Examíname, oh Dios, y conoce mi corazón\ldots{} Y ve si hay en mí camino de perversidad'' (Sal 139:23--24); ``Escudriñemos nuestros caminos'' (Lam 3:40); ``Examinaos a vosotros mismos'' (2~Co 13:5).
\item \textbf{Se endereza o cae.} Lo torcido que se somete es corregido: ``Reconócelo en todos tus caminos, Y él enderezará tus veredas'' (Pr 3:6); ``lo torcido se enderece'' (Is 40:4). Lo que se resiste termina por caer (Is 30:13).
\end{enumerate}

\begin{verso}{Proverbios 3:5–6}
Fíate de \SENOR{} de todo tu corazón, Y no te apoyes en tu propia prudencia. Reconócelo en todos tus caminos, Y él enderezará tus veredas.
\end{verso}

En Pablo, la medida llega a ser una Persona. La iglesia crece ``a la medida de la estatura de la plenitud de Cristo'' (\grc{εἰς μέτρον ἡλικίας τοῦ πληρώματος τοῦ Χριστοῦ}), ``para que ya no seamos niños fluctuantes, llevados por doquiera de todo viento de doctrina\ldots{} sino que siguiendo la verdad en amor, crezcamos en todo en aquel que es la cabeza, esto es, Cristo'' (Ef 4:13--15). Cristo es la plomada, y la deriva tiene nombre: ``todo viento de doctrina''.

%===================================================================
\section{Teología: el Dios que no miente}

La teología de la verdad de Dios no arranca de una definición, sino de una imposibilidad: la Escritura no se limita a decir que Dios no miente, sino que afirma que no puede mentir. El testimonio más claro procede de un testigo inesperado, Balaam, el profeta pagano a quien Balac contrató para maldecir a Israel y que, al abrir la boca, descubre que la palabra de Dios no se deja torcer:

\hebverso{לֹא אִישׁ אֵל וִיכַזֵּב וּבֶן־אָדָם וְיִתְנֶחָם}
\begin{verso}{Números 23:19}
Dios no es hombre, para que mienta, Ni hijo de hombre para que se arrepienta. El dijo, ¿y no hará? Habló, ¿y no lo ejecutará?
\end{verso}

\begin{filologia}{Nm 23:19}
\heb{וִיכַזֵּב} \emph{vijazzev} & \heb{כזב} & ``para que mienta'': el verbo de \emph{kazav}, la mentira que defrauda. \\
\heb{וְיִתְנֶחָם} \emph{veyitnejam} & \heb{נחם} & ``para que se arrepienta'': cambie de parecer. \\
\end{filologia}

Balaam argumenta por contraste: el hombre miente porque es hombre, y Dios no miente porque no lo es. La mentira no es, pues, un descuido que Dios evite con esfuerzo, sino algo ajeno a su ser. Samuel retoma el razonamiento cuando Saúl pierde el reino, y lo formula con la palabra que este capítulo ya ha identificado como el nombre de la mentira:

\hebverso{וְגַם נֵצַח יִשְׂרָאֵל לֹא יְשַׁקֵּר וְלֹא יִנָּחֵם}
\begin{verso}{1 Samuel 15:29}
Además, el que es la Gloria de Israel no mentirá, ni se arrepentirá, porque no es hombre para que se arrepienta.
\end{verso}

\begin{filologia}{1~S 15:29}
\heb{נֵצַח יִשְׂרָאֵל} \emph{nétsaj Yisra'el} & \heb{נצח} & ``la Gloria de Israel'': lo que permanece, lo perpetuo de Israel. \\
\heb{יְשַׁקֵּר} \emph{yeshaqqer} & \heb{שׁקר} & ``mentirá'': el verbo de \emph{shéqer}, la falsedad de los ídolos de Jer 10:14. \\
\end{filologia}

La Torá (Nm 23:19) y los Profetas anteriores (1~S 15:29) concuerdan, así, en que la veracidad pertenece al ser de Dios. Por la misma razón resulta inseparable de su fidelidad: una sola raíz, \emph{amán}, produce \emph{emet} y \emph{emuná}, y negar su palabra equivaldría a negarse a sí mismo. Pablo lo expresa en esos mismos términos: ``Si fuéremos infieles, él permanece fiel; él no puede negarse a sí mismo'' (2~Ti 2:13). Como observa Pink, esta cualidad ``es esencial para su Ser; sin ella no sería Dios. Para Dios, ser infiel, sería actuar en contra de su naturaleza, lo cual sería imposible'' (Pink, 2020, cap. 10). El Nuevo Testamento llega a resumirlo en un solo adjetivo, \grc{ἀψευδής}, ``que no miente'' (Tit 1:2), y la carta a los Hebreos asienta sobre esa imposibilidad toda la esperanza del creyente: ``es imposible que Dios mienta'' (He 6:18).

Si Dios no puede mentir, su palabra ha de verificarse en la historia, y la Escritura muestra que así sucede. Dios anunció a Abraham que su descendencia sería esclava en tierra ajena y volvería ``en la cuarta generación'' (Gn 15:13--16), y el Éxodo registra el cumplimiento con precisión de calendario: ``pasados los cuatrocientos treinta años, en el mismo día todas las huestes de \SENOR{} salieron de la tierra de Egipto'' (Éx 12:41). Dios anunció por Isaías que la virgen concebiría (Is 7:14), y Pablo ve llegar la hora señalada: ``cuando vino el cumplimiento del tiempo, Dios envió a su Hijo'' (Gá 4:4). Entre la promesa y su cumplimiento pueden transcurrir siglos, pero ninguna palabra se pierde en el camino. Por eso Pink concluye: ``Dios es veraz. La promesa de su Palabra es segura'' (Pink, 2020, cap. 10).

La consecuencia final es de orden práctico. Confesar que Dios no miente resulta fácil mientras no cuesta nada; la prueba llega cuando lo que se ve parece contradecir lo que Dios dijo. Pink lo advierte con precisión: ``una cosa es aceptar la fidelidad de Dios como una verdad divina y otra muy distinta es actuar de acuerdo con ella'' (Pink, 2020, cap. 10). Reaparece aquí la lógica de Isaías 7:9: creer es apoyarse en lo firme, y solo permanece quien se apoya. El que descansa en la palabra del Dios que no miente aprende a acallar la queja y a esperar, porque sabe que el ``Amén'' de Dios sigue en pie.

\precision{veracidad y fidelidad comparten la raíz \emph{amán}, pero este capítulo se ocupa de lo que concierne a la veracidad de Dios; la fidelidad tiene su propio estudio en esta serie (capítulo 15).}

%===================================================================
\section{Las esferas de la verdad en la creación}

Si Dios es la fuente de la verdad, lo verdadero alcanza a todas las disciplinas humanas, y cualquier pregunta sobre ciencia o sobre economía política es, en el fondo, una pregunta sobre cómo Dios estructuró la realidad. Cuando la ingeniería respeta la gravedad, la estructura se mantiene en pie porque se somete a la verdad material. Jeremías 10 atribuye ese orden al Dios verdadero, ``el que hizo la tierra con su poder, el que puso en orden el mundo con su saber, y extendió los cielos con su sabiduría'' (Jer 10:12), donde ``puso en orden'' traduce \heb{מֵכִין} (\emph{mejín}), del verbo \emph{kun}, ``establecer firme''.

\subsection{La ciencia descubre}
La ciencia no inventa la verdad, sino que la descubre. Cuando la física o las matemáticas hallan un patrón exacto en el universo, lo que encuentran es la verdad del diseño de Dios, porque las leyes que rigen la materia expresan en el plano físico realidades que él estableció. Pablo lo afirma al comienzo de Romanos:

\begin{verso}{Romanos 1:20}
Porque las cosas invisibles de él, su eterno poder y deidad, se hacen claramente visibles desde la creación del mundo, siendo entendidas por medio de las cosas hechas, de modo que no tienen excusa.
\end{verso}

Para ilustrarlo, el material de la serie recurre a las leyes de la termodinámica.

\subsection{Primera ley: la conservación de la energía}
\textbf{El postulado.} La energía no se crea ni se destruye; solo se transforma. El universo dispone de una cantidad finita de energía que va cambiando de estado.

\textbf{La correspondencia.} Crear de la nada es un acto exclusivo de Dios, y un acto ya concluido: ``Y acabó Dios en el día séptimo la obra que hizo'' (Gn 2:2). El universo no puede generar por sí mismo nueva existencia. La ley ilustra además el principio bíblico del aliento de vida: cuando una persona muere, lo que la animaba no se aniquila, sino que el cuerpo se degrada y el espíritu vuelve a su fuente.

\hebverso{וְיָשֹׁב הֶעָפָר עַל־הָאָרֶץ כְּשֶׁהָיָה וְהָרוּחַ תָּשׁוּב אֶל־הָאֱלֹהִים אֲשֶׁר נְתָנָהּ}
\begin{verso}{Eclesiastés 12:7}
y el polvo vuelva a la tierra, como era, y el espíritu vuelva a Dios que lo dio.
\end{verso}

\precision{la correspondencia es una analogía propuesta por el material de la serie. El \emph{rúaj} (\heb{רוּחַ}, ``aliento, espíritu'') no es energía física; lo que la ley ilustra es que la creación ni se crea a sí misma ni se aniquila por sí sola.}

\subsection{Segunda ley: la entropía}
\textbf{El postulado.} En un sistema cerrado, la energía utilizable se disipa con el tiempo y el desorden (la entropía) tiende a aumentar. Todo en el universo se encamina al desgaste, al enfriamiento y a la muerte, y nada se ordena a sí mismo de forma espontánea.

\textbf{La correspondencia.} El material de la serie ve aquí la demostración física de la Caída y de la maldición del pecado: el universo fue creado bueno, pero tras la rebelión del hombre la creación entera quedó sometida a un principio de degradación y muerte. Lo que la ciencia llama ``aumento de entropía'', la Biblia lo llama ``esclavitud de corrupción''.

\begin{verso}{Romanos 8:20–21}
Porque la creación fue sujetada a vanidad, no por su propia voluntad, sino por causa del que la sujetó en esperanza; porque también la creación misma será libertada de la esclavitud de corrupción, a la libertad gloriosa de los hijos de Dios.
\end{verso}
\begin{verso}{Salmo 102:25–27 [102:26–28]}
Desde el principio tú fundaste la tierra, Y los cielos son obra de tus manos. Ellos perecerán, mas tú permanecerás; Y todos ellos como una vestidura se envejecerán; Como un vestido los mudarás, y serán mudados; Pero tú eres el mismo, Y tus años no se acabarán.
\end{verso}

\begin{filologia}{Ro 8:20--21}
\grc{ματαιότης} \emph{mataiótēs} & --- & ``vanidad'': frustración, lo que no alcanza su fin. En la LXX traduce el \emph{hével} de Eclesiastés. \\
\grc{δουλεία τῆς φθορᾶς} & --- & ``esclavitud de la corrupción'': \grc{φθορά} es descomposición, deterioro. \\
\end{filologia}

El salmista contempla un universo que se gasta como ropa vieja mientras el Creador permanece, y Hebreos aplicará estas palabras a Cristo (He 1:10--12).

\precision{identificar la segunda ley con la Caída es la lectura del material de la serie, y es discutida. Otros intérpretes entienden que la segunda ley operaba ya en la creación buena (el sol, por ejemplo, alumbra gastando su energía) y leen Ro 8:20 como la frustración que siguió al pecado, sin identificarla con una ley física. Lo que el texto sí afirma es que la creación actual está sometida a corrupción y aguarda su liberación.}

\subsection{Gobierno, política y economía}
La verdad alcanza también al orden social, cuyo punto de partida es el mandato cultural primario de Génesis 1:28:

\hebverso{פְּרוּ וּרְבוּ וּמִלְאוּ אֶת־הָאָרֶץ וְכִבְשֻׁהָ וּרְדוּ}
\begin{verso}{Génesis 1:28}
Y los bendijo Dios, y les dijo: Fructificad y multiplicaos; llenad la tierra, y sojuzgadla, y señoread en los peces del mar, en las aves de los cielos, y en todas las bestias que se mueven sobre la tierra.
\end{verso}

\begin{filologia}{Gn 1:28}
\heb{וְכִבְשֻׁהָ} \emph{vejivshuha} & \heb{כבשׁ} & ``y sojuzgadla'': poner bajo orden lo que está sin cultivar. \\
\heb{וּרְדוּ} \emph{uredú} & \heb{רדה} & ``y señoread'': gobernar como rey, con la autoridad delegada del Creador. \\
\end{filologia}

Lejos de autorizar una explotación caprichosa, este mandato funda el orden social y material, y es anterior a la Caída (Gn 2:15). El pastor Miguel Núñez lo expresa así: ``Cada pregunta sobre política y economía es una pregunta sobre lo que significa ser hecho a la imagen de Dios y tener dominio sobre la Tierra. Dios le dijo al hombre: `sojuzga la tierra y domínala'. Bueno, para sojuzgar la tierra y dominarla habría que desarrollar algo que se llama ciencias políticas y ciencias de la economía''.

\precision{esta cita de Núñez procede del material de la serie (una predicación) y todavía falta su fuente exacta. En su libro \emph{Renueva tu mente}, el autor trata el mandato de Gn 1:28 como fundamento del trabajo humano (Núñez, 2020, cap. 9).}

Medidas con la plomada de la verdad, estas disciplinas quedan definidas en su esencia:
\begin{itemize}
\item \textbf{Ciencias políticas.} No se reducen a buscar y retener el poder; consisten en diseñar sistemas de gobierno, leyes y estructuras sociales que reflejen la justicia, el orden y el carácter del Creador, con el fin de gobernar y proteger a sociedades formadas por portadores de la imagen de Dios. Toda tiranía, corrupción o legislación que atenta contra ese diseño (como la opresión del vulnerable o la alteración de la familia) opera bajo un sistema de mentira.
\item \textbf{Economía.} Es la mayordomía fiel y eficiente de los recursos de la creación, que organiza el trabajo humano, la producción, el intercambio y la distribución de bienes para el florecimiento de la sociedad. Una economía fundada en la verdad somete el mundo material al orden divino y enfrenta la escasez con la creatividad humana, a diferencia de los sistemas corrompidos por la avaricia o el control opresivo.
\end{itemize}

\subsection{El fundamento de la moral}
Sin la verdad absoluta de Dios no pueden existir un ser humano ni una sociedad morales. Cuando se retira la plomada, el hombre acomoda la moral a su conveniencia y se vuelve capaz de ignorar el sufrimiento ajeno. En 1996, en el Everest, un grupo de alpinistas japoneses pasó junto a escaladores que agonizaban sin ayudarlos; uno de ellos lo justificó diciendo que ``por encima de los 8000 metros no es un lugar donde la gente pueda permitirse la moralidad'' (Krakauer, 1997). Privada de la verdad, la moral cede ante la conveniencia.

\begin{verso}{Salmo 12:1–2 [12:2–3]}
Salva, oh \SENOR{}, porque se acabaron los piadosos; Porque han desaparecido los fieles de entre los hijos de los hombres. Habla mentira cada uno con su prójimo; Hablan con labios lisonjeros, y con doblez de corazón.
\end{verso}

Los ``fieles'' son los \heb{אֱמוּנִים} (\emph{emunim}), los que viven en \emph{emuná}; y cuando ellos desaparecen, con ellos desaparece la verdad entre prójimos.

\precision{la frase se atribuye a un miembro de la expedición japonesa de Fukuoka; conviene verificar la página en la edición del libro de Krakauer que se utilice.}

%===================================================================
\section{La verdad frente a la percepción humana}

\subsection{“Todo hombre es mentiroso”}
Desde la Caída, el ser humano distorsiona a menudo la realidad, unas veces para justificar sus acciones y otras porque interpreta los hechos desde sus propios límites.

\hebverso{אֲנִי אָמַרְתִּי בְחָפְזִי כָּל־הָאָדָם כֹּזֵב}
\begin{verso}{Salmo 116:11}
Y dije en mi apresuramiento: Todo hombre es mentiroso.
\end{verso}

Cuando Pablo cita este salmo en griego, le antepone la otra mitad de la verdad:

\begin{verso}{Romanos 3:3–4}
¿Pues qué, si algunos de ellos han sido incrédulos? ¿Su incredulidad habrá hecho nula la fidelidad de Dios? De ninguna manera; antes bien sea Dios veraz, y todo hombre mentiroso.
\end{verso}

\begin{center}\small\renewcommand{\arraystretch}{1.2}
\begin{tabular}{@{}L{2.6cm}L{11.8cm}@{}}
\toprule
Sal 116:11 TM & \heb{כָּל־הָאָדָם כֹּזֵב} \\
Sal 115:2 LXX & \grc{πᾶς ἄνθρωπος ψεύστης} \\
Ro 3:4 & \grc{γινέσθω δὲ ὁ θεὸς ἀληθής, πᾶς δὲ ἄνθρωπος ψεύστης} \\
\bottomrule
\end{tabular}\par
{\footnotesize\color{gris}\textit{Tabla 6.} Sal 116:11 en el texto masorético, la Septuaginta y Romanos.}
\end{center}

\subsection{Llamar a lo malo bueno}
\begin{verso}{Isaías 5:20}
¡Ay de los que a lo malo dicen bueno, y a lo bueno malo; que hacen de la luz tinieblas, y de las tinieblas luz; que ponen lo amargo por dulce, y lo dulce por amargo!
\end{verso}

El ``¡ay!'' (\heb{הוֹי}, \emph{hoy}) es el lamento de los funerales, y el oráculo recuerda que cambiar el nombre de las cosas no cambia las cosas. En tiempos de Jeremías, el profeta Hananías anunció lo que el pueblo deseaba oír, la paz en dos años (Jer 28:3), y Jeremías le replicó:

\begin{verso}{Jeremías 28:15}
Entonces dijo el profeta Jeremías al profeta Hananías: Ahora oye, Hananías: \SENOR{} no te envió, y tú has hecho confiar en mentira a este pueblo.
\end{verso}

``Hiciste confiar'' traduce \heb{הִבְטַחְתָּ} (\emph{hivtajta}): Hananías dio seguridad al pueblo, pero una seguridad apoyada en \emph{shéqer}. La mentira cómoda se cree sin esfuerzo, mientras que la verdad se reconoce porque se cumple (Dt 18:22; Jer 28:9, 17).

\subsection{Los sentimientos y la realidad biológica}
El material de la serie critica la idea moderna según la cual ``la verdad es lo que yo siento''. La verdad no descansa en los sentimientos humanos: si lo que alguien siente no corresponde a su realidad biológica o genética, ese sentimiento no es verdad. El caso más discutido en la actualidad es el de la identidad sexual, cuya realidad señalan el diseño genético (cromosomas XX o XY), el sistema hormonal y la estructura ósea y cerebral: ``Varón y hembra los creó'' (Gn 1:27).

Núñez desarrolla el punto en su libro sobre la cosmovisión: ``Por primera vez en la historia podemos ver que una verdad, en este caso biológica, es definida en base a un sentimiento. Algo totalmente absurdo\ldots{} Explicar la verdad tal como es, podría ayudar más a estas personas a aceptar su realidad que mutilar su cuerpo para adaptarlo a una irrealidad de la mente'' (Núñez, 2020, cap. 10). En una predicación lo formuló de este modo: ``No es tu cuerpo el que tenemos que ajustar con cirugías, es tu mente la que tenemos que ajustar con la verdad y llevarla a la verticalidad de la verdad de Dios''.

\precision{la segunda cita de Núñez procede del material de la serie (una predicación) y todavía falta su fuente exacta; la primera se toma de la edición de 2020 de \emph{Renueva tu mente}.}

\subsection{Documentación clínica y legal}
Alterar la biología para acomodarla a una percepción tiene consecuencias graves, como muestran tres casos recientes que llegaron a los tribunales de Estados Unidos:

\begin{enumerate}
\item \textbf{Allan Josephson.} Psiquiatra infantil y jefe de la División de Psiquiatría y Psicología Infantil y Adolescente de la Universidad de Louisville, fue degradado tras afirmar, en 2017, que la disforia de género en niños ``no puede abordarse plenamente con fármacos y cirugía'', y en 2019 no se le renovó el contrato. En 2024 el Tribunal de Apelaciones del Sexto Circuito resolvió que su demanda debía ir a juicio, y en abril de 2025 la universidad acordó pagarle 1,6 millones de dólares.
\item \textbf{Fox Varian.} En 2019, a los 16 años y con un diagnóstico de autismo, se le practicó una mastectomía doble. El 30 de enero de 2026 un jurado de White Plains (Nueva York) condenó al psicólogo y al cirujano por mala praxis, en particular por no haber obtenido un consentimiento informado adecuado, y le otorgó 2 millones de dólares (1,6 por daño y 0,4 por gastos médicos futuros).
\item \textbf{Kaley Chiles.} En el caso de esta consejera cristiana de Colorado, la Corte Suprema de Estados Unidos resolvió el 31 de marzo de 2026, por 8 votos contra 1, que el Estado no puede aplicar su prohibición de las terapias de conversión para menores a la conversación terapéutica de Chiles, porque la ley discrimina según el punto de vista y debe pasar el escrutinio más estricto de la Primera Enmienda (\emph{Chiles v. Salazar}, 2026).
\end{enumerate}

\precision{conviene presentar los casos según lo que efectivamente resolvieron. En el caso Josephson no hubo sentencia sobre el fondo, sino un acuerdo previo al juicio. En el caso Varian, el jurado se pronunció sobre el estándar de cuidado médico y el consentimiento informado. El caso Chiles es un fallo sobre libertad de expresión: la ley de Colorado prohibía a los consejeros ayudar a un menor a cambiar su orientación o identidad, y la Corte decidió que el Estado no puede silenciar uno de los lados de una cuestión moral y psicológica disputada. Fuentes: Spectrum News (21 de abril de 2025); ABA Journal (4 de febrero de 2026) y Alliance Defending Freedom; opinión de la Corte en \emph{Chiles v. Salazar}.}

%===================================================================
\section{Un encuentro en la historia de Israel: la viuda de Sarepta}

En tiempos de Acab, Elías anunció una sequía: ``Vive \SENOR{} Dios de Israel, en cuya presencia estoy, que no habrá lluvia ni rocío en estos años, sino por mi palabra'' (1~R 17:1). El anuncio desafiaba a Baal, el dios cananeo de la lluvia y la tormenta. Dios envió luego a Elías a Sarepta, en Sidón, la tierra de Jezabel, donde una viuda le dio de comer y ``la harina de la tinaja no escaseó, ni el aceite de la vasija menguó, conforme a la palabra que \SENOR{} había dicho por Elías'' (17:16). Más tarde murió el hijo de aquella mujer:

\begin{verso}{1 Reyes 17:21–24}
Y se tendió sobre el niño tres veces, y clamó a \SENOR{} y dijo: \SENOR{} Dios mío, te ruego que hagas volver el alma de este niño a él. Y \SENOR{} oyó la voz de Elías, y el alma del niño volvió a él, y revivió. Tomando luego Elías al niño, lo trajo del aposento a la casa, y lo dio a su madre, y le dijo Elías: Mira, tu hijo vive. Entonces la mujer dijo a Elías: Ahora conozco que tú eres varón de Dios, y que la palabra de \SENOR{} es verdad en tu boca.
\end{verso}

\hebverso{עַתָּה זֶה יָדַעְתִּי כִּי אִישׁ אֱלֹהִים אָתָּה וּדְבַר־יְהוָה בְּפִיךָ אֱמֶת}

\begin{filologia}{1 R 17:24}
\heb{עַתָּה זֶה יָדַעְתִּי} \emph{attá ze yadati} & \heb{ידע} & ``ahora sí sé'': un conocimiento nuevo, por experiencia. \\
\heb{וּדְבַר־יְהוָה בְּפִיךָ אֱמֶת} & \heb{אמן} & ``y la palabra del \SENOR{} en tu boca es \emph{emet}'': firme, confiable, verdadera. \\
\heb{רְאִי חַי בְּנֵךְ} \emph{re'í jay benej} & \heb{חיה} & ``mira, tu hijo vive'' (17:23). \\
\end{filologia}

La mujer ya llamaba a Elías ``varón de Dios'' (17:18), pero ahora \emph{conoce}: la palabra es \emph{emet} porque ha devuelto la vida. En la tierra de Baal, el dios que no pudo dar lluvia, el \SENOR{} devolvió un hijo a su madre. El relato prepara así la confesión del Carmelo, ``¡\SENOR{} es el Dios, \SENOR{} es el Dios!'' (1~R 18:39), y reproduce el contraste de Jer 10:10: el Dios verdadero es el Dios vivo.

Jesús evoca a esta viuda en la sinagoga de Nazaret: ``muchas viudas había en Israel en los días de Elías\ldots{} pero a ninguna de ellas fue enviado Elías, sino a una mujer viuda en Sarepta de Sidón'' (Lc 4:25--26). Más adelante, en Naín, resucita al hijo único de otra viuda, y Lucas cierra el relato con las mismas palabras de la Septuaginta: \grc{καὶ ἔδωκεν αὐτὸν τῇ μητρὶ αὐτοῦ}, ``y lo dio a su madre'' (Lc 7:15 = 1~R 17:23 LXX).

\precision{el eco verbal entre 1~R 17:23 (LXX) y Lc 7:15 consta en el texto griego; la conexión entre la confesión de la viuda y Jer 10:10, en cambio, es una lectura propia de esta serie.}

%===================================================================
\section{La verdad en el judaísmo}\label{sec:judaismo-verdad}

\subsection{“El sello del Santo es \emph{emet}”}
Al comentar la marca que el ángel pone en la frente de los justos en Ezequiel 9, el Talmud de Babilonia recoge un dicho de Rabí Janiná:

\hebverso{חוֹתָמוֹ שֶׁל הַקָּדוֹשׁ בָּרוּךְ הוּא אֱמֶת}

``El sello del Santo, bendito sea, es \emph{emet}'' (\emph{Shabat} 55a). El Talmud de Jerusalén conserva la misma tradición y la desarrolla (\emph{Sanedrín} 1:1): a la pregunta ``¿qué es el sello del Santo?'' responde ``\emph{emet}'', y Rabí Bun aclara su sentido con las palabras de Jer 10:10, \heb{שֶׁהוּא אֱלֹהִים חַיִּים וּמֶלֶךְ עוֹלָם}, ``que él es Dios vivo y Rey eterno''. Resh Laquish añade que álef es la primera letra del alfabeto, mem la del medio y tav la última, porque Dios es el primero, que no recibió de otro, y el último, que no entregará a otro (Is 41:4; cf. 44:6).

\precision{interpretación rabínica. Texto según Sefaria (Talmud de Babilonia, edición William Davidson; Talmud de Jerusalén, edición de Guggenheimer). La lectura de las tres letras es un juego interpretativo, no una etimología.}

\subsection{La verdad tiene pies anchos}
En otra página, el Talmud compara las palabras \emph{shéqer} (\heb{שקר}) y \emph{emet} (\heb{אמת}). Las letras de \emph{shéqer} están juntas en el alfabeto (\heb{ש ק ר}) y las de \emph{emet}, separadas, porque la mentira abunda y la verdad escasea. Observa además que las letras de \emph{shéqer} se apoyan sobre un solo punto, mientras que las de \emph{emet} descansan sobre una base ancha:

\hebverso{קוּשְׁטָא קָאֵי שִׁיקְרָא לָא קָאֵי}

``La verdad se mantiene en pie; la mentira no se mantiene'' (\emph{Shabat} 104a). Los sabios formulan aquí, a su manera, la plomada del material de la serie y la misma idea de Pr 12:19.

\precision{interpretación rabínica sobre la forma de las letras hebreas cuadradas, según Sefaria.}

\subsection{La verdad arrojada a la tierra}
El midrash \emph{Bereshit Rabá} (8:5) relata, en nombre de Rabí Simón, que cuando Dios se disponía a crear al hombre los ángeles se dividieron. La Misericordia dijo ``que sea creado'', porque hará actos de bondad; la Verdad replicó ``que no sea creado'', \heb{שֶׁכֻּלּוֹ שְׁקָרִים}, ``porque todo él es mentiras''; y también la Justicia y la Paz entraron en la disputa. El Santo, entonces, \heb{נָטַל אֱמֶת וְהִשְׁלִיכוֹ לָאָרֶץ}, ``tomó la verdad y la arrojó a la tierra'', como dice Dn 8:12 (``echó por tierra la verdad''). Los sabios cierran el relato con el Salmo:

\hebverso{אֱמֶת מֵאֶרֶץ תִּצְמָח וְצֶדֶק מִשָּׁמַיִם נִשְׁקָף}
\begin{verso}{Salmo 85:10–11 [85:11–12]}
La misericordia y la verdad se encontraron; La justicia y la paz se besaron. La verdad brotará de la tierra, Y la justicia mirará desde los cielos.
\end{verso}

\precision{interpretación rabínica (midrash). La lectura cristiana que ve en Cristo ``la verdad que brota de la tierra'' (cf. Jn 1:14) no pertenece al midrash.}

\subsection{La verdad en la oración y en la Misná}
\begin{itemize}
\item \textbf{Después del \emph{Shemá}.} El \emph{Shemá} concluye con ``Yo soy el \SENOR{} vuestro Dios'' (Nm 15:41), y el Talmud enseña que a continuación se dice \emph{emet}, ``verdad'', apoyándose en Jer 10:10: \heb{וַיהוָה אֱלֹהִים אֱמֶת} (\emph{Berajot} 14b). En la oración diaria, el Nombre y la verdad van siempre juntos.
\item \textbf{\emph{Emet veyatsiv}.} La bendición matinal que sigue al \emph{Shemá} comienza con \heb{אֱמֶת וְיַצִּיב}, ``verdadero y firme'', y la vespertina, con \heb{אֱמֶת וֶאֱמוּנָה}, ``verdad y fidelidad'' (Elbogen, 1993).
\item \textbf{\emph{Avot} 1:18.} Rabán Shimón ben Gamliel enseña: \heb{עַל שְׁלֹשָׁה דְבָרִים הָעוֹלָם עוֹמֵד}, ``sobre tres cosas se sostiene el mundo: la justicia, la verdad y la paz'', y lo funda en Zac 8:16.
\item \textbf{Maimónides.} En el primer capítulo de la \emph{Mishné Torá} (\emph{Yesodé haTorá} 1:4) cita Jer 10:10 y lo comenta así: \heb{הוּא לְבַדּוֹ הָאֱמֶת וְאֵין לְאַחֵר אֱמֶת כַּאֲמִתָּתוֹ}, ``él solo es la verdad, y ningún otro tiene verdad como la suya''.
\end{itemize}

\begin{verso}{Zacarías 8:16}
Estas son las cosas que habéis de hacer: Hablad verdad cada cual con su prójimo; juzgad según la verdad y lo conducente a la paz en vuestras puertas.
\end{verso}

\precision{interpretaciones y usos litúrgicos judíos, según Sefaria. Conviene notar que tres de las cuatro fuentes se apoyan en el texto ancla de este capítulo, Jer 10:10.}

\subsection{Qumrán: el espíritu de verdad}
La \emph{Regla de la Comunidad} de Qumrán (1QS) enseña que Dios puso en el hombre dos espíritus para andar en ellos hasta el tiempo de su visitación: \heb{רוּחוֹת הָאֱמֶת וְהָעָוֶל}, ``los espíritus de la verdad y de la perversidad'' (1QS 3:18--19). ``Los nacidos de la verdad brotan de una fuente de luz, y los nacidos de la perversidad, de un manantial de tinieblas'' (1QS 3:19). Al final, Dios ``rociará sobre él el espíritu de verdad como aguas de purificación'' (1QS 4:21), con un lenguaje que recuerda Ez 36:25--27; y quien ingresa en la comunidad se compromete ``a practicar la verdad, la justicia y el derecho en la tierra'' (1QS 1:5), en un eco de Jer 9:24.

El Evangelio de Juan emplea un vocabulario muy cercano: ``el Espíritu de verdad, al cual el mundo no puede recibir'' (Jn 14:17) y ``el que practica la verdad viene a la luz'' (Jn 3:21).

\precision{traducción según García Martínez (1992) y Vermes (2004); conviene cotejar las imágenes de 1QS en la Leon Levy Dead Sea Scrolls Digital Library. Los estudios del Nuevo Testamento señalan los paralelos con Juan, que revelan un lenguaje judío común y no una dependencia directa. Sobre Qumrán y el texto de Jeremías 10, véase la precisión textual sobre Jeremías 10 en la sección~\ref{sec:jer10}.}

%===================================================================
\section{La verdad en Cristo, leída con Pablo}\label{sec:cristo-verdad}

\subsection{Cambiaron la verdad por la mentira}
Pablo abre su exposición del evangelio con el mismo diagnóstico que Jeremías 10:

\begin{verso}{Romanos 1:18, 25}
Porque la ira de Dios se revela desde el cielo contra toda impiedad e injusticia de los hombres que detienen con injusticia la verdad\ldots{} ya que cambiaron la verdad de Dios por la mentira, honrando y dando culto a las criaturas antes que al Creador, el cual es bendito por los siglos. Amén.
\end{verso}

\begin{filologia}{Ro 1:18, 25}
\grc{τὴν ἀλήθειαν\ldots{} κατεχόντων} & \grc{κατέχω} & ``que sujetan, retienen la verdad'': la verdad está a la vista, pero la reprimen. \\
\grc{μετήλλαξαν} \emph{metéllaxan} & \grc{ἀλλάσσω} & ``trocaron'': cambiaron una cosa por otra. \\
\grc{ἐν τῷ ψεύδει} & \grc{ψεῦδος} & ``por \emph{la} mentira'', con artículo: el ídolo. La LXX de Jer 10:14 dice de los fundidores \grc{ψευδῆ ἐχώνευσαν}, ``fundieron mentiras''. \\
\grc{ἀμήν} & --- & Pablo cierra la denuncia con un ``Amén'': frente a la mentira, la confesión de lo firme. \\
\end{filologia}

Para Pablo, como para Jeremías, ``la mentira'' es ante todo el ídolo; y la verdad de Dios es Dios mismo, conocido por medio de lo creado (Ro 1:20) y trocado por una imagen.

\precision{leer ``la mentira'' de Ro 1:25 como la idolatría es habitual entre los comentaristas (Moo, 1996); el vínculo con Jer 10:14, en cambio, es una lectura propia de esta serie.}

\subsection{“Sea Dios veraz”}
En Romanos 3 Pablo sale al paso de una objeción: si Israel fue infiel, ¿no queda anulada la fidelidad de Dios?

\begin{verso}{Romanos 3:3–4}
¿Pues qué, si algunos de ellos han sido incrédulos? ¿Su incredulidad habrá hecho nula la fidelidad de Dios? De ninguna manera; antes bien sea Dios veraz, y todo hombre mentiroso; como está escrito: Para que seas justificado en tus palabras, Y venzas cuando fueres juzgado.
\end{verso}

La respuesta une las dos caras de la raíz \emph{amán}. La ``incredulidad'' (\grc{ἀπιστία}) de algunos no anula ``la fidelidad de Dios'' (\grc{τὴν πίστιν τοῦ θεοῦ}), y Dios es ``veraz'' (\grc{ἀληθής}) porque cumple lo que prometió aunque Israel falle; en 3:7, Pablo llama a eso mismo ``la verdad de Dios'' (\grc{ἡ ἀλήθεια τοῦ θεοῦ}). Es Nm 23:19 releído en Romanos: Dios no es hombre para que mienta.

\subsection{Cristo, siervo de la verdad de Dios}
\begin{verso}{Romanos 15:8–9}
Pues os digo, que Cristo Jesús vino a ser siervo de la circuncisión para mostrar la verdad de Dios, para confirmar las promesas hechas a los padres, y para que los gentiles glorifiquen a Dios por su misericordia.
\end{verso}

\begin{filologia}{Ro 15:8}
\grc{ὑπὲρ ἀληθείας θεοῦ} & \grc{ἀλήθεια} & ``a favor de la verdad de Dios'': para mostrar que Dios es fiel. \\
\grc{εἰς τὸ βεβαιῶσαι} & \grc{βεβαιόω} & ``para confirmar, hacer firmes'' las promesas: casi una traducción de \emph{amán}. \\
\end{filologia}

Miqueas había pedido: ``Cumplirás la verdad a Jacob, y a Abraham la misericordia, que juraste a nuestros padres desde tiempos antiguos'' (Mi 7:20). Pablo responde que Cristo trajo la verdad (\emph{emet}) a Israel y la misericordia (\emph{jésed}) a las naciones, de modo que Jesús se convierte en la prueba histórica de que Dios es \emph{Elohim emet}.

\subsection{“En él Sí, y en él Amén”}
\begin{verso}{2 Corintios 1:18–20}
Mas, como Dios es fiel, nuestra palabra a vosotros no es Sí y No. Porque el Hijo de Dios, Jesucristo, que entre vosotros ha sido predicado por nosotros, por mí, Silvano y Timoteo, no ha sido Sí y No; mas ha sido Sí en él; porque todas las promesas de Dios son en él Sí, y en él Amén, por medio de nosotros, para la gloria de Dios.
\end{verso}

\grc{ὅσαι γὰρ ἐπαγγελίαι θεοῦ, ἐν αὐτῷ τὸ ναί· διὸ καὶ δι᾽ αὐτοῦ τὸ ἀμήν}. El ``Sí'' es la respuesta de Dios a sus propias promesas en Cristo; el ``Amén'', la respuesta de la iglesia, que lo confiesa ``para la gloria de Dios''. Isaías había llamado a Dios \emph{Elohé amén} (Is 65:16), y el Apocalipsis llamará a Cristo ``el Amén, el testigo fiel y verdadero'' (Ap 3:14).

\begin{center}
\begin{tikzpicture}[font=\small,node distance=0.9cm]
 \node[draw=azul,line width=0.8pt,fill=papiro,rounded corners=3pt,text width=3.4cm,align=center,inner sep=6pt] (a) {\textbf{Is 65:16}\\ \heb{אֱלֹהֵי אָמֵן}\\ ``el Dios del Amén''};
 \node[draw=azul,line width=0.8pt,fill=papiro,rounded corners=3pt,text width=3.4cm,align=center,inner sep=6pt,right=of a] (b) {\textbf{2 Co 1:20}\\ \grc{δι᾽ αὐτοῦ τὸ ἀμήν}\\ ``en él Amén''};
 \node[draw=dorado,line width=1.6pt,fill=azul,text=papiro,rounded corners=3pt,text width=3.4cm,align=center,inner sep=6pt,right=of b] (c) {\textbf{Ap 3:14}\\ \grc{ὁ ἀμήν}\\ ``el Amén, el testigo fiel y verdadero''};
 \draw[-{Latex[length=3mm]},dorado,line width=1.4pt] (a)--(b);
 \draw[-{Latex[length=3mm]},dorado,line width=1.4pt] (b)--(c);
\end{tikzpicture}\par
{\footnotesize\color{gris}\textit{Figura 3.} Del Dios del Amén al Amén encarnado. \textit{Nota.} Elaboración propia.}
\end{center}

\subsection{El Dios vivo y verdadero}
\begin{verso}{1 Tesalonicenses 1:9–10}
porque ellos mismos cuentan de nosotros la manera en que nos recibisteis, y cómo os convertisteis de los ídolos a Dios, para servir al Dios vivo y verdadero, y esperar de los cielos a su Hijo, al cual resucitó de los muertos, a Jesús, quien nos libra de la ira venidera.
\end{verso}

Primera Tesalonicenses es, probablemente, la carta más antigua de Pablo, y al resumir en ella su predicación a los gentiles recurre al lenguaje de Jeremías 10: dejar los ídolos para servir \grc{θεῷ ζῶντι καὶ ἀληθινῷ}, ``al Dios vivo y verdadero'' (cf. \heb{אֱלֹהִים אֱמֶת}\ldots{} \heb{אֱלֹהִים חַיִּים}, Jer 10:10). La prueba de que Dios está vivo es que resucitó a su Hijo, algo que en Sarepta solo se había anticipado.

\precision{Jer 10:10 falta en la Septuaginta, y Pablo emplea una fórmula judía corriente (cf. Is 65:16 LXX, \grc{τὸν θεὸν τὸν ἀληθινόν}); el eco de Jeremías 10 es, por tanto, probable, pero no seguro.}

\subsection{La verdad está en Jesús}
El Antiguo Testamento preparó a Israel al revelarle un Dios que es firme y que exige firmeza. La revelación culminante del Nuevo Testamento consiste en mostrar que la verdad no es un simple concepto ni una idea filosófica, sino una Persona.

\begin{verso}{Efesios 4:20–21}
Mas vosotros no habéis aprendido así a Cristo, si en verdad le habéis oído, y habéis sido por él enseñados, conforme a la verdad que está en Jesús.
\end{verso}

La verdad está ``en Jesús'' (\grc{καθώς ἐστιν ἀλήθεια ἐν τῷ Ἰησοῦ}). Juan lo expresa con las palabras de Éxodo 34:6 (\emph{rav jésed ve'emet}): el Verbo ``fue hecho carne\ldots{} lleno de gracia y de verdad'' (Jn 1:14); ``la ley por medio de Moisés fue dada, pero la gracia y la verdad vinieron por medio de Jesucristo'' (Jn 1:17). Y Jesús lo afirma de sí mismo:

\begin{verso}{Juan 14:6}
Jesús le dijo: Yo soy el camino, y la verdad, y la vida; nadie viene al Padre, sino por mí.
\end{verso}

Al declarar ``Yo soy\ldots{} la verdad'', Jesucristo habla de su propia naturaleza encarnada. Por eso Pablo advierte: ``Mirad que nadie os engañe por medio de filosofías y huecas sutilezas\ldots{} y no según Cristo. Porque en él habita corporalmente toda la plenitud de la Deidad'' (Col 2:8--9). Lo que no está a plomo con Cristo tampoco lo está con la realidad.

\precision{Juan no es Pablo; se lo cita porque Pablo dice lo mismo con otras palabras. La autoría de Efesios y Colosenses es discutida, pero Ro 15:8 y 2~Co 1:20 afirman lo mismo en cartas indiscutidas.}

\subsection{La piedra de cimiento estable}\label{sec:piedra}
En Isaías 28, Dios pone en Sion una piedra probada y mide con cordel y plomada, y el que cree (\emph{hama'amín}) no queda en vergüenza. Pablo cita dos veces ese oráculo y lo aplica a Cristo:

\begin{verso}{Romanos 9:33}
como está escrito: He aquí pongo en Sion piedra de tropiezo y roca de caída; Y el que creyere en él, no será avergonzado.
\end{verso}
\begin{verso}{Romanos 10:11}
Pues la Escritura dice: Todo aquel que en él creyere, no será avergonzado.
\end{verso}

\begin{center}\small\renewcommand{\arraystretch}{1.2}
\begin{tabular}{@{}L{2.4cm}L{12.0cm}@{}}
\toprule
Is 28:16 TM & \heb{מוּסָד מוּסָּד הַמַּאֲמִין לֹא יָחִישׁ}: ``cimiento bien cimentado; el que cree no se apresurará'' \\
Is 28:16 LXX & \grc{ὁ πιστεύων ἐπ᾽ αὐτῷ οὐ μὴ καταισχυνθῇ}: ``el que cree en él no será avergonzado'' \\
Ro 10:11 & \grc{πᾶς ὁ πιστεύων ἐπ᾽ αὐτῷ οὐ καταισχυνθήσεται}: Pablo agrega ``todo aquel'' \\
\bottomrule
\end{tabular}\par
{\footnotesize\color{gris}\textit{Tabla 7.} Is 28:16 en el TM, la LXX y Romanos.}
\end{center}

Con ello la cadena queda completa: la raíz \emph{amán} es sostener, \emph{emet} es lo que sostiene y \emph{he'emín} es apoyarse en ello; Cristo, finalmente, es la piedra firme sobre la que se apoya el que cree.

\subsection{Mentir profana a una Persona}
Si la verdad es una Persona, las implicaciones son eternas. Quien miente o se alía con la falsedad no deshonra solamente un concepto ético abstracto, sino que profana a una Persona, a Jesús, y se pone del lado del padre de mentira.

\begin{verso}{Juan 8:44}
Vosotros sois de vuestro padre el diablo, y los deseos de vuestro padre queréis hacer. El ha sido homicida desde el principio, y no ha permanecido en la verdad, porque no hay verdad en él. Cuando habla mentira, de suyo habla; porque es mentiroso, y padre de mentira.
\end{verso}

Pablo extrae de ahí la consecuencia para la iglesia:

\begin{verso}{Efesios 4:25}
Por lo cual, desechando la mentira, hablad verdad cada uno con su prójimo; porque somos miembros los unos de los otros.
\end{verso}

\grc{λαλεῖτε ἀλήθειαν ἕκαστος μετὰ τοῦ πλησίον αὐτοῦ} reproduce casi palabra por palabra Zac 8:16 en la Septuaginta (\grc{λαλεῖτε ἀλήθειαν ἕκαστος πρὸς τὸν πλησίον αὐτοῦ}). Lo nuevo es la razón: ``somos miembros los unos de los otros''. Mentir al hermano es mentir dentro del cuerpo de Cristo.

\subsection{La verdad que hace libres}
\begin{verso}{Juan 8:31–32}
Si vosotros permaneciereis en mi palabra, seréis verdaderamente mis discípulos; y conoceréis la verdad, y la verdad os hará libres.
\end{verso}

La libertad nace de \emph{permanecer} (\grc{μένω}) en su palabra, y de nuevo aparece la firmeza. Pablo describe su ministerio como ``manifestación de la verdad'' y lo funda en la gloria contemplada en Cristo: ``Dios, que mandó que de las tinieblas resplandeciese la luz, es el que resplandeció en nuestros corazones, para iluminación del conocimiento de la gloria de Dios en la faz de Jesucristo'' (2~Co 4:2, 6). Solo en sujeción a Cristo, la Verdad encarnada y la plomada absoluta del universo, se hallan la salvación y la libertad.

%===================================================================
\section{Aplicación: un atributo comunicable}

\subsection{Hablar verdad en el corazón}
\hebverso{הוֹלֵךְ תָּמִים וּפֹעֵל צֶדֶק וְדֹבֵר אֱמֶת בִּלְבָבוֹ}
\begin{verso}{Salmo 15:1–2}
\SENOR{}, ¿quién habitará en tu tabernáculo? ¿Quién morará en tu monte santo? El que anda en integridad y hace justicia, Y habla verdad en su corazón.
\end{verso}

El salmo dice ``habla verdad \emph{en} su corazón'' (\heb{בִּלְבָבוֹ}): la verdad comienza en el interior, antes de llegar a los labios, como confirma otro salmo: ``He aquí, tú amas la verdad en lo íntimo'' (Sal 51:6 [51:8]). En cuanto fuente, la verdad pertenece solo a Dios; en cuanto veracidad, se comunica a quien vive delante de él.

\subsection{Siguiendo la verdad en amor}
Pablo pide a la iglesia que crezca ``siguiendo la verdad en amor'' (\grc{ἀληθεύοντες δὲ ἐν ἀγάπῃ}, Ef 4:15). El verbo \grc{ἀληθεύω} significa a la vez ``decir la verdad'' y ``vivir en la verdad''. La verdad sin amor hiere, y el amor sin verdad engaña. Zacarías ya lo había formulado como mandato para el tiempo de la restauración: ``Amad, pues, la verdad y la paz'' (\heb{וְהָאֱמֶת וְהַשָּׁלוֹם אֱהָבוּ}, Zac 8:19).

\subsection{Cuatro respuestas}
\begin{enumerate}
\item \textbf{Dejar los ídolos.} Servir al Dios vivo y verdadero, y no a lo que nosotros mismos fabricamos y clavamos para que no se caiga (Jer 10:4, 10; 1~Ts 1:9).
\item \textbf{Ajustarse a la plomada.} Permitir que la Palabra mida la mente y los sentimientos: ``Santifícalos en tu verdad; tu palabra es verdad'' (Jn 17:17); ``transformaos por medio de la renovación de vuestro entendimiento'' (Ro 12:2).
\item \textbf{Hablar verdad.} Hacerlo con el prójimo y en amor, porque formamos un solo cuerpo (Ef 4:15, 25; Zac 8:16).
\item \textbf{Amar la verdad.} Recibir ``el amor de la verdad'' en lugar de la mentira cómoda (2~Ts 2:10; Zac 8:19).
\end{enumerate}

\begin{verso}{2 Tesalonicenses 2:10}
y con todo engaño de iniquidad para los que se pierden, por cuanto no recibieron el amor de la verdad para ser salvos.
\end{verso}

El material de la serie deja abierta una pregunta que alcanza a cada lector: si estamos dispuestos a ajustar nuestra mente, nuestras emociones y nuestra vida a la plomada de la verdad de Dios, o si esperamos que la verdad se incline ante lo que sentimos.

%===================================================================
\section{Síntesis}

\begin{tcolorbox}[colback=papiro,colframe=azul,boxrule=1pt,arc=3pt]
Frente a unos ídolos clavados para que no se caigan, Jeremías confiesa: ``Mas \SENOR{} es el Dios verdadero'' (\emph{Elohim emet}, Jer 10:10). La verdad de Dios es firmeza (raíz \emph{amán}), correspondencia con la realidad que él creó y fidelidad a lo que promete. Como plomada de la ciencia, la política y la moral, juzga también lo que sentimos. Israel la llamó el ``sello'' de Dios, y Pablo la vio confirmada en Cristo, en quien todas las promesas son ``Sí'' y ``Amén'' (2~Co 1:20).
\end{tcolorbox}

\begin{center}\small\renewcommand{\arraystretch}{1.3}
\begin{tabular}{@{}L{3.0cm}L{11.4cm}@{}}
\toprule
\textbf{Mirada} & \textbf{Aporte al atributo} \\ \midrule
Hebrea & \emph{emet}, \emph{emuná}, \emph{amén}; Jeremías 10 (el Dios vivo frente al ídolo \emph{shéqer}); la plomada (Am 7; Is 28); Pr 12:19; la viuda de Sarepta. \\
Judía & El sello de Dios es \emph{emet}; \emph{shéqer} sobre un pie; \emph{emet} después del \emph{Shemá}; \emph{Avot} 1:18; Maimónides; el espíritu de verdad en 1QS. \\
Teológica & El Dios que no puede mentir (Nm 23:19; 1~S 15:29; Tit 1:2; He 6:18): veracidad esencial a su ser, verificada en la historia (Gn 15; Éx 12:41; Gá 4:4), para vivirla en la prueba. \\
Paulina & Ro 1:25; 3:4; 15:8; 2~Co 1:20; 1~Ts 1:9; Ef 4:21, 25; Cristo, la piedra firme (Ro 9:33; 10:11). \\
\bottomrule
\end{tabular}\end{center}

\begin{verso}{Jeremías 9:24 [9:23]}
Mas alábese en esto el que se hubiere de alabar: en entenderme y conocerme, que yo soy \SENOR{}, que hago misericordia, juicio y justicia en la tierra; porque estas cosas quiero, dice \SENOR{}.
\end{verso}

El verso que articula esta obra invita a conocer a Dios, y el capítulo siguiente de Jeremías revela quién es: el Dios verdadero, vivo y Rey eterno. Conocerlo equivale a apoyarse en lo único que no se cae.

%===================================================================
\section*{Compendio bíblico}
\begin{center}\small\renewcommand{\arraystretch}{1.12}
\begin{tabular}{@{}L{2.2cm}L{5.0cm}|L{2.2cm}L{5.0cm}@{}}
\toprule
Gn 1:28 & ``llenad la tierra, y sojuzgadla'' & Jn 1:14 & ``lleno de gracia y de verdad'' \\
Sal 15:1--2 & ``Y habla verdad en su corazón'' & Jn 3:21 & ``el que practica la verdad viene a la luz'' \\
Sal 26:3 & ``Y ando en tu verdad'' & Jn 4:23--24 & ``en espíritu y en verdad'' \\
Sal 31:5 & ``oh \SENOR{}, Dios de verdad'' & Jn 8:32 & ``la verdad os hará libres'' \\
Sal 43:3 & ``Envía tu luz y tu verdad'' & Jn 8:44 & ``es mentiroso, y padre de mentira'' \\
Sal 58:3 & ``hablando mentira desde que nacieron'' & Jn 14:6 & ``Yo soy el camino, y la verdad'' \\
Sal 86:11 & ``caminaré yo en tu verdad'' & Jn 14:17; 15:26 & ``el Espíritu de verdad'' \\
Sal 116:11 & ``Todo hombre es mentiroso'' & Jn 16:13 & ``os guiará a toda la verdad'' \\
Sal 119:30 & ``Escogí el camino de la verdad'' & Ro 1:25 & ``cambiaron la verdad de Dios por la mentira'' \\
Sal 119:89 & ``Permanece tu palabra en los cielos'' & 1~Co 3:18 & ``Nadie se engañe a sí mismo'' \\
Sal 119:160 & ``La suma de tu palabra es verdad'' & 2~Co 1:20 & ``en él Sí, y en él Amén'' \\
Pr 6:16--19 & ``la lengua mentirosa\ldots{} El testigo falso'' & Gá 6:7 & ``Dios no puede ser burlado'' \\
Pr 19:5 & ``el que habla mentiras no escapará'' & Ef 5:6 & ``Nadie os engañe con palabras vanas'' \\
Ec 12:7 & ``el espíritu vuelva a Dios que lo dio'' & Col 2:8 & ``nadie os engañe por medio de filosofías'' \\
Sal 102:25--26 & ``como una vestidura se envejecerán'' & 2~Ti 2:15 & ``que usa bien la palabra de verdad'' \\
Is 5:20 & ``a lo malo dicen bueno'' & He 6:18 & ``es imposible que Dios mienta'' \\
Is 65:16 & ``en el Dios de verdad se bendecirá'' & Stg 1:16 & ``Amados hermanos míos, no erréis'' \\
Jer 5:1 & ``si hay alguno\ldots{} que busque verdad'' & 1~P 1:22 & ``por la obediencia a la verdad'' \\
Jer 10:10 & ``Mas \SENOR{} es el Dios verdadero'' & 1~Jn 1:8; 3:7 & ``la verdad no está en nosotros''; ``nadie os engañe'' \\
Zac 8:19 & ``Amad, pues, la verdad y la paz'' & 3~Jn 3, 8 & ``tu verdad''; ``cooperemos con la verdad'' \\
 & & Ap 16:7; 19:2 & ``tus juicios son verdaderos y justos'' \\
 & & Ap 19:9, 11 & ``palabras verdaderas''; ``Fiel y Verdadero'' \\
\bottomrule
\end{tabular}\par
{\footnotesize\color{gris}\textit{Tabla 8.} Compendio bíblico de la verdad. \textit{Nota.} Textos del material de la serie más el texto ancla, en la RVR1960.}
\end{center}

%===================================================================
\section*{Glosario del capítulo}
\begin{description}\setlength{\itemsep}{2pt}
\item[\emph{Alétheia}] (\grc{ἀλήθεια}). Verdad; en la LXX traduce \emph{emet}.
\item[\emph{Alēthinós}] (\grc{ἀληθινός}). Genuino, verdadero, real (1~Ts 1:9; Jn 15:1).
\item[\emph{Amán}] (\heb{אָמַן}). Sostener, cargar; ser firme; en hifil, creer.
\item[\emph{Amén}] (\heb{אָמֵן}, \grc{ἀμήν}). ``Así es'', ``firme''; ``el Dios del Amén'' (Is 65:16); ``el Amén'' (Ap 3:14).
\item[\emph{Anaj}] (\heb{אֲנָךְ}). Plomada (o estaño) (Am 7:7--8).
\item[\emph{Apseudés}] (\grc{ἀψευδής}). Que no miente (Tit 1:2).
\item[\emph{Elohim emet}] (\heb{אֱלֹהִים אֱמֶת}). ``Dios verdad'' (Jer 10:10).
\item[\emph{Dípsijos}] (\grc{δίψυχος}). ``De doble alma'', con doblez (Stg 1:8).
\item[\emph{Emet}] (\heb{אֱמֶת}). Verdad, firmeza, fidelidad.
\item[\emph{Emet veyatsiv}] (\heb{אֱמֶת וְיַצִּיב}). ``Verdadero y firme'': bendición matinal después del \emph{Shemá}.
\item[\emph{Emuná}] (\heb{אֱמוּנָה}). Fidelidad, firmeza; fe.
\item[\emph{He'emín}] (\heb{הֶאֱמִין}). Creer: apoyarse en lo firme (Is 7:9).
\item[\emph{Hével}] (\heb{הֶבֶל}). Vanidad, vapor (Jer 10:3).
\item[\emph{Iqqesh}] (\heb{עִקֵּשׁ}). Torcido, desviado (Dt 32:5).
\item[\emph{Jotam}] (\heb{חוֹתָם}). Sello: ``el sello del Santo es \emph{emet}''.
\item[\emph{Kazav}] (\heb{כָּזָב}). Mentira; el verbo: mentir (Sal 116:11; Nm 23:19).
\item[\emph{Mishquélet}] (\heb{מִשְׁקֶלֶת}). Nivel, plomada (Is 28:17; 2~R 21:13).
\item[\emph{Ne'emán}] (\heb{נֶאֱמָן}). Firme, fiel, estable.
\item[\emph{Omnam}] (\heb{אָמְנָם}). Verdaderamente.
\item[\emph{Phōs}] (\grc{φῶς}). Luz: ``Dios es luz'' (1~Jn 1:5).
\item[\emph{Pseûdos}] (\grc{ψεῦδος}). Mentira (Ro 1:25; Ef 4:25).
\item[\emph{Qeshot}] (arameo \heb{קְשֹׁט}). Verdad (Dn 4:37).
\item[\emph{Shéqer}] (\heb{שֶׁקֶר}). Mentira, falsedad, engaño (Jer 10:14).
\item[\emph{Yashar}] (\heb{יָשָׁר}). Recto, derecho, a plomo (Dt 32:4; Ec 7:29).
\end{description}

